\documentclass[11pt]{article}

\usepackage[T1]{fontenc}
\usepackage{amsmath,amssymb,amsthm,mathtools}
\usepackage[margin=1.1in]{geometry}
\usepackage[hidelinks]{hyperref}

\title{Rigidity of Fourier Multipliers for Universal Euler Convective
Cancellations}
\author{Rael Nogueira Pires}
\date{}

\newtheorem{theorem}{Theorem}[section]
\newtheorem{corollary}[theorem]{Corollary}
\newtheorem{lemma}[theorem]{Lemma}
\newtheorem{proposition}[theorem]{Proposition}
\theoremstyle{definition}
\newtheorem{definition}[theorem]{Definition}
\theoremstyle{remark}
\newtheorem{remark}[theorem]{Remark}

\theoremstyle{plain}
\newtheorem*{thmA}{Theorem A}
\newtheorem*{thmAzero}{Theorem A$^{0}$}
\newtheorem*{corB}{Corollary B}
\newtheorem*{corC}{Corollary C}
\newtheorem*{thmD}{Theorem D}
\newtheorem*{corE}{Corollary E}

\DeclareMathOperator{\curl}{curl}
\DeclareMathOperator{\dv}{div}
\DeclareMathOperator{\Sym}{Sym}
\DeclareMathOperator{\rank}{rank}
\DeclareMathOperator{\nullity}{nullity}
\DeclareMathOperator{\spn}{span}
\DeclareMathOperator{\range}{range}
\DeclareMathOperator{\cond}{cond}
\DeclareMathOperator{\dist}{dist}
\DeclareMathOperator{\sgn}{sgn}
\DeclareMathOperator{\tr}{tr}
\newcommand{\T}{\mathbb{T}}
\newcommand{\Z}{\mathbb{Z}}
\newcommand{\R}{\mathbb{R}}
\newcommand{\C}{\mathbb{C}}
\newcommand{\Q}{\mathbb{Q}}
\newcommand{\Qform}{\mathcal{Q}}
\newcommand{\Dform}{\mathcal{D}}
\newcommand{\Ck}{C_k}
\newcommand{\kperp}{k_{\C}^{\perp}}
\newcommand{\gsym}{\nabla_{\mathrm{sym}}}

\begin{document}

\emergencystretch=2em

\maketitle

\begin{abstract}
We prove an operator-level rigidity theorem for universal Euler convective
cancellations on the three-torus. Let
$R(k):\kperp\to\C^{3}$, $k\in\Z^{3}\setminus\{0\}$, be an arbitrary
reality-compatible Fourier symbol, with no boundedness, locality, isotropy,
self-adjointness, or range assumption. If
$\langle Rw,(w\cdot\nabla)w\rangle_{L^{2}}=0$ for every real, mean-zero,
divergence-free trigonometric polynomial $w$, then, and only then,
\[
R=A+B\curl+\curl B,
\qquad A,B\in\operatorname{Sym}(3,\R).
\]
Thus universal cancellation forces an arbitrary symbol to be the restriction
of a constant-coefficient self-adjoint operator of order at most one. If the
output is also divergence-free, the family reduces to
$R=cI+d\curl$. The zero-mode block is invisible and remains arbitrary when
constant fields are admitted. The proof combines phase-polarized triads,
analytic propagation through $\Z^{3}$, and exact finite rank certificates.
We also obtain sharp finite test counts under explicit ordering hypotheses,
identify the isosceles degeneracy caused by Leray projection, and pull the
classification back to universal strain--vorticity cancellations. The
certificates are reproduced by exact-arithmetic programs; the Lean~4
formalization and its compiler-trusted finite-rank dependency are documented
separately. No application to PDE regularity or existence is claimed.
\end{abstract}

\tableofcontents

\section{Introduction}
\label{sec:intro}

\subsection{The classification problem}
\label{subsec:problem}

Work on the three-torus $\T^{3}=(\R/2\pi\Z)^{3}$ with normalized Haar
measure, and call a field \emph{admissible} if it is a real, mean-zero,
divergence-free trigonometric polynomial. The incompressible Euler equation
transports such a field by the quadratic convective nonlinearity
$(w\cdot\nabla)w$. Conservation of kinetic energy and of helicity are the
statements that the two particular operators $I$ and $\curl$ are orthogonal to
that nonlinearity,
\begin{equation}
\label{eq:intro-energy-helicity}
\int_{\T^{3}}w\cdot(w\cdot\nabla)w\,dx=0,
\qquad
\int_{\T^{3}}\curl w\cdot(w\cdot\nabla)w\,dx=0,
\end{equation}
for every admissible $w$. We ask for the converse. Let
\begin{equation}
\label{eq:intro-symbol}
R(k):\kperp\longrightarrow\C^{3},
\qquad k\in\Z^{3}\setminus\{0\},
\end{equation}
be an arbitrary complex-linear symbol, assumed only to be compatible with
reality, and define the cubic functional
\begin{equation}
\label{eq:intro-Q}
\Qform_{R}(w)
:=\int_{\T^{3}}Rw\cdot(w\cdot\nabla)w\,dx .
\end{equation}
\emph{Which symbols satisfy $\Qform_{R}(w)=0$ for every admissible $w$?}

The hypotheses omit the assumptions that usually make such a question
tractable. The symbol may depend on $k$ arbitrarily; it need not be bounded,
homogeneous, isotropic, local, Hermitian, or transverse in its range. Each
fiber is an independent complex $3\times2$ matrix, so the symbol space is
infinite-dimensional and has no prescribed structure. Nevertheless, the
universal cancellation condition forces the symbol into a
twelve-dimensional real family obtained by restricting constant-coefficient,
formally self-adjoint operators of order at most one, and a two-dimensional
family once incompressibility of the output is imposed. Allowing longitudinal
output enlarges the cancellation kernel; imposing the solenoidal range
condition collapses it to the energy--helicity family. Within the
reality-compatible,
translation-invariant Fourier-multiplier class with solenoidal output considered
here, the two operators in \eqref{eq:intro-energy-helicity} exhaust the
possibilities.

\subsection{Main results}
\label{subsec:main-results}

Write $\Ck v=ik\times v$ for the curl symbol, so that $\Ck$ is Hermitian on
$\kperp$ and $\Ck^{2}=|k|^{2}I$ there. The main classification is the
following; it is restated and proved as Theorem~A in
Sections~\ref{sec:results} and~\ref{sec:propagation}.

\begin{quote}
\textbf{Theorem A.} A complex-linear, reality-compatible symbol
$R(k):\kperp\to\C^{3}$ satisfies $\Qform_{R}(w)=0$ for every admissible $w$ if
and only if there are unique constant real symmetric matrices $A,B$ with
\[
R(k)v=Av+B\Ck v+\Ck Bv,
\qquad v\in\kperp,\quad k\ne0 .
\]
\end{quote}

The result has four useful consequences.

\emph{Hermiticity is a conclusion.} The pairing in \eqref{eq:intro-Q} is the
real $L^{2}$ pairing; in Fourier variables it is the complex-bilinear extension
of the Euclidean dot product, pairing the coefficients at $k$ and $-k$. No
adjoint is inserted anywhere, and the hypothesis is not that a quadratic form
is conserved. A non-Hermitian symbol has a genuine skew sector which
\emph{does} contribute to $\Qform_{R}$, and the phase polarization of
Lemma~\ref{lem:phase} is exactly what detects it. Every classified symbol is nevertheless the restriction to $\kperp$ of a
formally self-adjoint constant-coefficient operator on the full vector-field
space. In the solenoidal class this is sharp:
Corollary~B below concludes $R=cI+d\curl$ with no Hermitian assumption made
anywhere.

\emph{Order and locality are conclusions.} The hypothesis allows nonlocal
symbols of arbitrary growth. The conclusion is a differential operator of
order at most one with constant coefficients.

\emph{The longitudinal directions are cancellations, not invariants.} The
functional \eqref{eq:intro-Q} is cubic in $w$ and pairs $R$ against the
quadratic convection. When $R$ is self-adjoint on the solenoidal phase space,
$\Qform_{R}$ is the Euler derivative of the quadratic functional
$\tfrac12\langle w,Rw\rangle$. For a general $R:\kperp\to\C^{3}$ it is only a
convective cancellation: longitudinal output can couple to the pressure, so
the ten extra directions in the twelve-dimensional family must not be
advertised as new Euler invariants. We return to this point in
Section~\ref{subsec:not-invariants}.

\emph{The zero mode is invisible, not rigid.} If the mean-zero hypothesis is
dropped, Theorem~A$^{0}$ shows that the classification is unchanged on every
nonzero mode while $R(0)$ remains an arbitrary real $3\times3$ matrix. No
field evaluation can see it, because periodic convection has zero spatial
mean. The total kernel dimension is then $12+9=21$.

Imposing incompressibility of the output collapses the family:

\begin{quote}
\textbf{Corollary B.} If in addition $R(k)\kperp\subset\kperp$ for every
$k\ne0$, then $\Qform_{R}\equiv0$ holds if and only if $R(k)=cI+d\Ck$ for two
global real constants $c,d$. In particular every nonzero-mode fiber is
Hermitian.
\end{quote}

Because $\curl$ is invertible on mean-zero solenoidal periodic fields, the
whole classification transports to the strain--vorticity pairing. Writing
$\omega=\curl u$ and
$S=\gsym u$, and setting
$\Dform_{T}(u)=\langle\gsym(Tu),\omega\otimes\omega\rangle$, we obtain:

\begin{quote}
\textbf{Corollary C.} $\Dform_{T}(u)=0$ for every admissible $u$ if and only if
$T=A\curl+B(-\Delta)+\curl B\curl$ with $A,B$ constant real symmetric; with
solenoidal output this reduces to $T=c\curl+d(-\Delta)$.
\end{quote}

The scalar case $T=-\Delta$ of the second formula gives the Laplacian identity
$\langle-\Delta S,\omega\otimes\omega\rangle=0$. Corollary~C says that within
the translation-invariant multiplier class this identity and the first-order
one are the only universal strain--vorticity cancellations with solenoidal
output, and it identifies precisely which longitudinal completions restore
cancellation for an arbitrary anisotropic second-order scalar operator
(Corollary~E, Section~\ref{subsec:anisotropic}).
The scalar identity is due to Miller~\cite{Miller2026}; the classification
here includes it as one member of the larger multiplier family.

Finally, the proof is quantitative. Under two explicit ordering hypotheses on
a finite symmetric mode set with $M$ unoriented classes, we prove the existence
of test families with exactly $8M-2$, $4M-2$ and $12M-12$ real field
evaluations that certify the classification in the complex solenoidal,
Hermitian solenoidal and unrestricted-output classes respectively, and we show
that no smaller family can work (Theorem~D). The machine-readable supplement
exports the selected seed rows, pivot columns, mode order and parameter
ordering used by the nonsingular-minor certificates. For a coordinate box of radius $N$ these have the
optimal $\Theta(N^{3})$ size rather than the $\Theta(N^{6})$ cost of scanning
a dense collection of triads.

\subsection{Method, and what is checked by machine}
\label{subsec:method}

The proof separates a finite exact calculation from analytic propagation over
the lattice.

Sufficiency (Section~\ref{sec:sufficiency}) is a short analytic computation.
For necessity, Lemma~\ref{lem:phase} shows that testing $\Qform_{R}$ on real
six-mode fields supported on a non-collinear triad $p+q+r=0$, with total phase
products $1$ and $i$, recovers the full complex triad coefficient. Two local
lemmas then describe what one triad sees at a fixed recipient mode. In the
unrestricted class, one triad annihilates eight of the twelve real parameters
of the recipient block and leaves a single complex image line
(Lemma~\ref{lem:rawline}); a second triad with a non-parallel difference
direction finishes the job. In the solenoidal class, a single triad with
parents of unequal length already determines all eight real parameters of the
recipient block, while equal parent lengths leave a four-dimensional residue
(Lemma~\ref{lem:projected}). This isosceles defect is created exactly by the
Leray projection and is invisible before it; Proposition~\ref{prop:projected}
computes the responsible singular value in closed form.

The induction that turns those local statements into a global theorem needs a
base case that no triad argument supplies, and this is where exact computation
enters. Two finite constraint matrices are built and their ranks are certified
over $\Q$:
\begin{center}
\begin{tabular}{lrrrrr}
\hline
Class & Seed modes & Parameters & Rows & Rank & Nullity\\
\hline
Hermitian solenoidal & $7$ & $28$ & $592$ & $26$ & $2$\\
Complex solenoidal & $7$ & $56$ & $720$ & $54$ & $2$\\
Complex unrestricted output & $13$ & $156$ & $680$ & $144$ & $12$\\
\hline
\end{tabular}
\end{center}
In each case the matching kernel is exhibited explicitly by the classified
generators, so rank and nullity together give an exact kernel \emph{equality},
not merely a bound (equations \eqref{eq:seed-kernel-sol} and
\eqref{eq:seed-kernel-raw}). Rank is certified from below by a nonsingular
minor evaluated modulo two primes, and from above by the explicit kernel. No
floating-point rank decision and no numerical tolerance enters at any point:
the programs use integers, exact rational arithmetic, or reductions modulo
primes. Two of the certificates are reproduced by algorithmically independent
implementations that reconstruct the convolution by recipient wavevector
rather than reusing the row builders.

The formal development in Lean~4 with Mathlib covers Theorem~A and
Corollary~B together with four supporting components: the arbitrary real
zero block in Theorem~A$^{0}$; the torus-integral realization of
Corollary~C and its extension to smooth solenoidal fields; the sharp scalar
evaluation counts underlying Theorem~D at
the finite-dimensional matrix level; and the local singular spectra,
conditioning bounds, bad family and stability modulo the exact kernel from
Section~\ref{sec:singular}. The source contains no admitted proofs. The two
finite seed rank certificates use \texttt{native\_decide}, which the
machine-readable axiom report exposes as seven compiled-code axioms, and the
twenty-six-row physical-to-matrix reconstruction for the Hermitian seed is an
external exact check rather than a kernel theorem. Section~\ref{sec:computer}
gives the result-by-result trust boundary and records what remains analytic.

\subsection{Relation to prior work}
\label{subsec:prior}

The result sits next to several classification literatures. We organize the
comparison by the object being classified rather than chronologically, since
the hypotheses differ far more than the conclusions do.

\emph{Rugged quadratic invariants and Galerkin truncations.} That energy and
helicity are well-established quadratic invariants of three-dimensional Euler
in Fourier-truncated settings. T.~D.~Lee's statistical mechanics of
Galerkin-truncated ideal flow~\cite{TDLee1952}, Kraichnan's
absolute-equilibrium theory and the Kraichnan--Montgomery
review~\cite{KraichnanMontgomery1980}, Jon Lee's triad-interaction
representation and isolating constants of motion~\cite{JLee1975a,JLee1975b},
his later study of mixing on constant energy--helicity
surfaces~\cite{JLee1982}, and Shebalin's monograph~\cite{Shebalin2002} form
the historical base. Shebalin states that three-dimensional Euler admits no
further bilinear integral invariants. None of these formulates an arbitrary
mode-dependent matrix symbol, retains longitudinal output, or treats the
non-Hermitian sector classified here.

\emph{Local conservation laws.} Serre classifies the first-order local
invariants of three-dimensional Euler~\cite{Serre1984}; Vulkov derives the
conservation laws and symmetric forms of the incompressible inviscid-fluid
equations~\cite{Vulkov1999}; and Cheviakov and Oberlack construct hierarchies
of conserved quantities from generalized Ertel
theorems~\cite{CheviakovOberlack2014}. These are differential-density
classifications under a local finite-order ansatz. The present hypothesis
admits arbitrary nonlocal modewise symbols from the start.

\emph{Hamiltonian and Casimir uniqueness.} Khesin and
Chekanov~\cite{KhesinChekanov1989}, Enciso, Peralta-Salas and Torres de
Lizaur~\cite{EncisoEtAl2016}, Khesin, Peralta-Salas and
Yang~\cite{KhesinEtAl2020}, and Dullin, Meiss and
Worthington~\cite{DullinEtAl2019} classify invariants under much stronger
geometric hypotheses --- invariance under volume-preserving diffeomorphisms,
or the Lie--Poisson structure in Fourier space. Helicity uniqueness there is a
statement about regular integral invariants of a group action. Our hypothesis
is convective orthogonality alone, and energy is included in the conclusion
rather than excluded by the symmetry requirement.

\emph{Helical and triad geometry.} Waleffe's helical decomposition makes the
energy--helicity constraints on a single triad transparent~\cite{Waleffe1992};
Moffatt exhibits further invariants of an isolated triad~\cite{Moffatt2014};
Rathmann and Ditlevsen study pseudo-invariants of restricted helical
interactions~\cite{RathmannDitlevsen2017}; Alexakis and Brachet analyze fluxes
in quasi-equilibrium~\cite{AlexakisBrachet2020}. Invariants of an isolated or
decimated triad system are strictly more numerous than universal ones, which
is exactly the distinction that propagation over the full lattice enforces
here.

\emph{Triad-network nullspaces.} Harper, Bustamante and
Nazarenko~\cite{HarperEtAl2013} reduce the quadratic invariants of discrete
clusters of weakly interacting waves to sparse matrix nullspaces. The broad
methodology of computing invariants as nullspaces of triad-interaction
matrices is therefore prior art, and we claim no novelty for it. What is
different here is the symbol class, the passage from a finite cluster to the
entire lattice, and the exact rank structure of that passage.

\emph{Kinematic strain--vorticity constraints.} Betchov's homogeneity
relation~\cite{Betchov1956} and the proof by Carbone and Wilczek that only two
such constraints exist for isotropic turbulence~\cite{CarboneWilczek2022} are
the closest kinematic antecedents of the pullback in
Section~\ref{sec:strainvorticity}. The scalar Laplacian
strain--vorticity identity is a direct antecedent of Corollary~C; our claim
concerns the classification and the anisotropic longitudinal completion.

\subsection{Novelty boundary and claims not made}
\label{subsec:boundary}

To the best of our knowledge, the existing local-conservation,
rugged-invariant, Hamiltonian/Casimir and helical-triad classifications do
not
state the all-symbol implication proved here, its automatic-Hermiticity
corollary, or the unrestricted longitudinal kernel. We therefore treat these
as candidate new results, subject to specialist and peer review. A targeted
prior-art audit located no exact duplicate, but a literature search cannot
establish historical priority, and equivalent statements could be embedded in
older sources indexed under different terminology.

The following standard facts are not claimed as new: conservation of energy and
helicity; their triad-by-triad manifestations; additional invariants of
isolated or decimated triad systems; the helical decomposition; classifications
of local conservation laws and of regular Casimirs; triad-network nullspace
methodology; and the scalar Laplacian strain--vorticity identity.

The following are explicitly \emph{not} proved anywhere in this paper.
\begin{enumerate}
\item No result on PDE existence, smoothness, or regularity is claimed. The
anisotropic completion of Section~\ref{subsec:anisotropic} is a structural identity; it identifies the
Hessian defect exactly, but supplies no new scale-critical, time-integrable
bound for it.
\item No closed coercive a priori estimate is proved.
\item No multiplier theorem on $\R^{3}$. The propagation argument uses the
lattice $\Z^{3}$ throughout.
\item No rigidity at the zero Fourier mode; that block is proved invisible and
remains arbitrary.
\item No boundedness, continuity or growth hypothesis is imposed on the
symbol. The classification itself implies
$\|R(k)\|=O_{A,B}(1+|k|)$, but no multiplier-extension or form-domain theorem on
completed function spaces is developed here. All tests are trigonometric
polynomials, so every sum is finite.
\item No classification of nonlinear, time-dependent, or
non-translation-invariant functionals.
\item Minimality of the seven-mode seed is proved only among subsets of the
coordinate unit cube, not among arbitrary six-mode lattice configurations.
\item Theorem~D applies only under its two explicit ordering hypotheses, not
to an arbitrary finite symmetric mode set.
\end{enumerate}

\subsection{Outline}
\label{subsec:outline}

Section~\ref{sec:setting} fixes the conventions --- frames, reality, the
bilinear pairing --- and proves the coordinate-invariance lemma that reconciles
the normalized frames used in the text with the integer frames used by the
certificate programs. Section~\ref{sec:results} states the complete theorem
package. Section~\ref{sec:sufficiency} proves sufficiency.
Section~\ref{sec:triad} develops the triad reduction and the two local lemmas.
Section~\ref{sec:seeds} constructs the exact finite seeds and states their
certified kernels. Section~\ref{sec:propagation} carries out the lattice
propagation and completes the proofs of Theorems~A, A$^{0}$, D and
Corollary~B. Section~\ref{sec:singular} computes the raw and projected
singular spectra and the conditioning of the explicit route.
Section~\ref{sec:strainvorticity} proves Corollary~C and the anisotropic
longitudinal completion. Section~\ref{sec:computer} documents the
computer-assisted components, their reproduction, and the formalization trust
boundary. Section~\ref{sec:discussion} collects limitations and open problems.
Appendix~\ref{app:certificates} tabulates the certificate data.

\section{Setting, notation and conventions}
\label{sec:setting}

\subsection{Fields and symbols}

Work on $\T^{3}=(\R/2\pi\Z)^{3}$ with normalized Haar measure. A
\emph{test field} is a real, mean-zero, divergence-free trigonometric
polynomial,
\begin{equation}
\label{eq:testfield}
w(x)=\sum_{k\in\Z^{3}\setminus\{0\}}\widehat w(k)e^{ik\cdot x},
\qquad
k\cdot\widehat w(k)=0,
\qquad
\widehat w(-k)=\overline{\widehat w(k)},
\end{equation}
with finite Fourier support. We call such a field \emph{admissible}. For each
$k\ne0$ set
\begin{equation}
\label{eq:curlsymbol}
\Ck v=ik\times v .
\end{equation}
On $\kperp$ the operator $\Ck$ is Hermitian and satisfies
\begin{equation}
\label{eq:curlsq}
\Ck^{2}=|k|^{2}I .
\end{equation}
A symbol $R$ is \emph{reality compatible} when
\begin{equation}
\label{eq:reality}
R(-k)\overline v=\overline{R(k)v},
\qquad v\in\kperp .
\end{equation}
Two symbol classes occur throughout,
\begin{equation}
\label{eq:classes}
\mathcal R_{\mathrm{raw}}
=\{R(k):\kperp\to\C^{3}\},
\qquad
\mathcal R_{\mathrm{sol}}
=\{R(k):\kperp\to\kperp\} ;
\end{equation}
the first allows longitudinal output, the second preserves
incompressibility. When the zero mode is admitted we use the convention
$0_{\C}^{\perp}=\C^{3}$, so that reality compatibility says exactly that
$R(0)$ is a real $3\times3$ matrix.

There is no boundedness or growth assumption on $R(k)$. Since every test
field is a trigonometric polynomial, all sums below are finite.

\subsection{The canonical transverse frame}

To make \eqref{eq:reality} an unambiguous matrix identity, fix a frame
convention. For $k\ne0$ let $\kappa(k)$ be the member of $\{k,-k\}$ whose
first nonzero coordinate is positive, let $j(\kappa)$ be the least index with
$\kappa\times e_{j}\ne0$, and set
\begin{equation}
\label{eq:frame}
f_{1}(\kappa)
=\frac{\kappa\times e_{j(\kappa)}}{|\kappa\times e_{j(\kappa)}|},
\qquad
f_{2}(\kappa)=\frac{\kappa}{|\kappa|}\times f_{1}(\kappa).
\end{equation}
For \emph{both} signed modes $k$ and $-k$ we use the same real orthonormal
frame
\begin{equation}
\label{eq:sameframe}
\mathcal F(k)=\mathcal F(-k)
=\bigl(f_{1}(\kappa(k)),f_{2}(\kappa(k))\bigr).
\end{equation}
Thus $(f_{1},f_{2},\kappa/|\kappa|)$ is positively oriented at the canonical
representative, and compatibility with reality --- rather than positive
orientation --- is what is retained at its negative. In these frames
\eqref{eq:reality} reads $R(-k)=\overline{R(k)}$.

\subsection{The pairing}
\label{subsec:pairing}

The pairing in \eqref{eq:intro-Q} is the real $L^{2}$ pairing of two real
fields. In Fourier coordinates it uses the complex-\emph{bilinear} extension
of the Euclidean dot product to pair the coefficient at $k$ with the
coefficient at $-k$. No adjoint and no conjugation is silently inserted.

This convention must be fixed before any row calculation, because it is what
makes the automatic-self-adjointness statement meaningful. Had the
sesquilinear pairing been used, the skew-Hermitian part of $R$ would be
annihilated by construction and Corollary~B would be a statement about a
quotient rather than about the symbol.

\subsection{Coordinate invariance of the certificates}

The exposition uses the normalized frames \eqref{eq:frame}; the certificate
programs use an integer, non-normalized transverse frame, which keeps all
arithmetic exact over $\Z$ or $\Q$. The following lemma reconciles them.

\begin{lemma}[coordinate and generated-row invariance]
\label{lem:coordinates}
Fix a finite geometric family of real triad-supported test fields. Replacing
only the transverse frames used to express the symbol blocks preserves the
rank, the nullity and the geometric kernel of the resulting constraint matrix.

Separately, fix finitely many triads and, on each of them, generate the full
row space using all tensor-product polarizations from a transverse basis at
each mode together with the two total phase classes $1$ and $i$. Replacing
those polarization bases preserves this complete generated row space, hence its
rank and nullity. No such statement is made for an arbitrary proper subset of
the coordinate-generated rows.
\end{lemma}

\begin{proof}
If $F_{k}$ and $G_{k}$ are two real bases of $k^{\perp}$, then
$G_{k}=F_{k}S_{k}$ with $S_{k}\in\mathrm{GL}_{2}(\R)$. For a fixed geometric
test family, changing the domain coordinates of a raw block --- or both domain
and range coordinates of a solenoidal block --- right-multiplies the complete
constraint matrix by an invertible block-diagonal matrix. This preserves rank
and nullity and merely re-expresses the same geometric kernel.

If the test polarizations are also replaced, trilinearity of
$(w,w,w)\mapsto\Qform_{R}(w)$ expresses the full tensor-product family in the
new bases as an invertible multilinear combination of the full family in the
old bases; Lemma~\ref{lem:phase} gives the same statement for the two real
phase rows. The two complete generated row spaces therefore agree after the
change of coordinates. An arbitrary selected subset of coordinate rows need
not be carried to an equally indexed subset, which is why no such claim is
included.
\end{proof}

The dense seed matrices of Section~\ref{sec:seeds} use the complete
phase--polarization families, so the second part of Lemma~\ref{lem:coordinates}
applies to them. Once particular minor rows have been selected, they are fixed
geometric test fields, and changing only the symbol coordinates falls under the
first part. The integer script convention therefore certifies the same
geometric kernels as the normalized convention used here. Since the changes of
coordinates need not be orthogonal, numerical singular values are \emph{not}
preserved; all singular values in Section~\ref{sec:singular} refer to the
normalized convention.

\subsection{Triad conventions}

A \emph{triad} is a relation $p+q+r=0$ among lattice modes, and it is
\emph{non-collinear} when $p\times q\ne0$. In propagation arguments the new
mode is called the \emph{target} and written $t$; its known \emph{parents} are
$q,r$ with
\begin{equation}
\label{eq:targetconv}
t=q+r,
\qquad
p=-t,
\qquad
p+q+r=0 .
\end{equation}
The local lemmas act at the \emph{recipient} $p=-t$ and therefore prove
$D(-t)=0$ first; reality compatibility then gives $D(t)=0$. Consequently
$D(k)=0$ if and only if $D(-k)=0$, so every set of known modes below is
automatically symmetric under $k\mapsto-k$, and we use this without further
comment.

The raw convection received at $p$ is
$N=(U\cdot r)V+(V\cdot q)U$; the Leray projection $P_{p^{\perp}}$ is applied
only in the solenoidal output class.

\section{The theorem package}
\label{sec:results}

\subsection{Unrestricted-output classification}
\label{subsec:thmA}

\begin{thmA}[unrestricted-output rigidity]
\label{thm:A}
Let $R(k):\kperp\to\C^{3}$, $k\ne0$, be a complex-linear, reality-compatible
symbol. Then
\begin{equation}
\label{eq:universal}
\Qform_{R}(w)=0
\quad\text{for every admissible }w
\end{equation}
holds if and only if there are unique constant real symmetric matrices $A,B$
with
\begin{equation}
\label{eq:generator}
R(k)v=Av+B\Ck v+\Ck Bv,
\qquad v\in\kperp,\quad k\ne0 .
\end{equation}
\end{thmA}

Sufficiency is proved in Section~\ref{sec:sufficiency}, necessity in
Section~\ref{subsec:rawprop}. Uniqueness is the following self-contained
statement, which uses only the signed coordinate-axis modes and is independent
of every finite-rank certificate.

\begin{proposition}[uniqueness of the two matrix coefficients]
\label{prop:uniqueness}
If $A,B\in\Sym(3,\R)$ satisfy
$Av+B\Ck v+\Ck Bv=0$ for every $k\ne0$ and every $v\in\kperp$, then $A=B=0$.
Consequently the pair $(A,B)$ in Theorem~A is unique.
\end{proposition}

\begin{proof}
Evaluate the identity at $k$ and at $-k$, recalling $C_{-k}=-\Ck$, and add and
subtract. For every $v\perp k$ this gives
\[
Av=0,
\qquad
B\Ck v+\Ck Bv=0 .
\]
Taking $k=e_{1}$ and $k=e_{2}$ in the first identity shows that $A$ annihilates
all three coordinate vectors, so $A=0$.

For the second identity write $\Ck=iK_{k}$ with $K_{k}v=k\times v$ and divide
by $i$. With $B=(b_{mn})$ symmetric, the four coordinate evaluations are
\[
\begin{array}{c|c}
(k,v)&BK_{k}v+K_{k}Bv\\ \hline
(e_{1},e_{2})&(b_{13},0,b_{22}+b_{33})^{T}\\
(e_{1},e_{3})&(-b_{12},-b_{22}-b_{33},0)^{T}\\
(e_{2},e_{1})&(0,-b_{23},-b_{11}-b_{33})^{T}\\
(e_{3},e_{1})&(0,b_{11}+b_{22},b_{23})^{T}.
\end{array}
\]
Setting these four vectors to zero gives
$b_{13}=b_{12}=b_{23}=0$ and
$b_{22}+b_{33}=b_{11}+b_{33}=b_{11}+b_{22}=0$, whence all three diagonal
entries vanish as well. Thus $B=0$.
\end{proof}

\begin{thmAzero}[all-mode completion]
\label{thm:Azero}
Let $R(k):\kperp\to\C^{3}$, $k\in\Z^{3}$, be complex-linear and reality
compatible, with $0_{\C}^{\perp}=\C^{3}$. Then
\begin{equation}
\label{eq:universal-allmode}
\Qform_{R}(w)=0
\quad\text{for every real divergence-free trigonometric polynomial }w
\end{equation}
holds if and only if there are unique $A,B\in\Sym(3,\R)$ such that
\eqref{eq:generator} holds for every $k\ne0$, while $R(0)$ is an arbitrary
real $3\times3$ matrix.
\end{thmAzero}

\begin{proof}
Necessity on the nonzero modes follows by restricting
\eqref{eq:universal-allmode} to mean-zero fields and applying Theorem~A;
reality compatibility at $k=0$ is exactly reality of the matrix entries.

For sufficiency write $w=a+u$ with $a\in\R^{3}$ constant and $u$ mean zero.
Periodicity and $\dv u=0$ give
\begin{equation}
\label{eq:zeromean}
\int_{\T^{3}}(a\cdot\nabla)u\,dx=0,
\qquad
\int_{\T^{3}}(u\cdot\nabla)u\,dx=0,
\end{equation}
so the constant vector $R(0)a$ is orthogonal to the complete convective term.
On the nonzero modes let $S=A+B\curl+\curl B$. This is a constant-coefficient
self-adjoint operator and it commutes with the directional derivative
$D_{a}=a\cdot\nabla$, whence
\begin{equation}
\label{eq:commute}
\langle Su,D_{a}u\rangle
=\langle u,D_{a}Su\rangle
=-\langle D_{a}u,Su\rangle=0 .
\end{equation}
The remaining term is $\Qform_{S}(u)=0$ by Theorem~A, or directly by
Section~\ref{sec:sufficiency}. Thus \eqref{eq:universal-allmode} holds for
every $w$, with no restriction on $R(0)$.
\end{proof}

Three consequences deserve emphasis. The universal kernel has real dimension
$12$ on the mean-zero phase space and dimension $21$ in the all-mode
formulation --- twelve rigid nonzero-mode parameters and nine invisible
zero-mode ones. A completely arbitrary, potentially nonlocal symbol is forced
to have constant coefficients and differential order at most one. And admitting
constant velocities as tests neither creates new nonzero-mode symbols nor
constrains the zero-mode block.

\subsection{Solenoidal classification and automatic self-adjointness}
\label{subsec:corB}

\begin{lemma}[scalarity from lattice directions]
\label{lem:scalarity}
Let $L\in\Sym(3,\R)$. If $Lk\parallel k$ for every $k\in\Z^{3}\setminus\{0\}$,
then $L=cI$ for some $c\in\R$.
\end{lemma}

\begin{proof}
Applying the hypothesis to $e_{1},e_{2},e_{3}$ makes the coordinate vectors
eigenvectors, so $L=\mathrm{diag}(\lambda_{1},\lambda_{2},\lambda_{3})$.
Applying it to $e_{i}+e_{j}$ gives $\lambda_{i}=\lambda_{j}$ for each $i\ne j$.
\end{proof}

\begin{corB}[solenoidal rigidity]
\label{cor:B}
Under the hypotheses of Theorem~A, assume in addition the range condition
\begin{equation}
\label{eq:range}
R(k)\kperp\subset\kperp
\qquad(k\ne0).
\end{equation}
Then \eqref{eq:universal} holds if and only if
\begin{equation}
\label{eq:scalarform}
R(k)=cI+d\Ck
\end{equation}
for two global real constants $c,d$. In particular the cubic identity forces
every nonzero-mode block $R(k)$ to be Hermitian.
\end{corB}

\begin{proof}
Sufficiency is contained in Theorem~A, since $cI+d\Ck$ is the case
$A=cI$, $B=(d/2)I$ of \eqref{eq:generator} and preserves $\kperp$ by
\eqref{eq:curlsq}. For necessity, apply Theorem~A to obtain
\eqref{eq:generator} and impose \eqref{eq:range}. Let $v\perp k$. Since
$\Ck Bv=ik\times Bv$ is orthogonal to $k$, the range condition at $k$ says
\[
k\cdot Av+k\cdot B\Ck v=0 ,
\]
and at $-k$, using $C_{-k}=-\Ck$ and the same frame \eqref{eq:sameframe}, it
says $k\cdot Av-k\cdot B\Ck v=0$. Adding and subtracting gives
\[
k\cdot Av=0,
\qquad
k\cdot B\Ck v=0
\qquad\text{for all }v\in\kperp .
\]
By \eqref{eq:curlsq} the map $\Ck:\kperp\to\kperp$ is onto, so the second
identity says $k\cdot Bw=0$ for every $w\in\kperp$. Symmetry of $A$ and $B$
now converts both statements into
\[
Ak\parallel k,
\qquad
Bk\parallel k
\]
for every lattice direction $k$. Lemma~\ref{lem:scalarity} forces $A=cI$ and
$B=bI$, and \eqref{eq:generator} becomes $R=cI+2b\Ck$.
\end{proof}

In the all-mode formulation of Theorem~A$^{0}$ the same argument gives
$R(k)=cI+d\Ck$ for every $k\ne0$, while $R(0)$ remains an arbitrary real
$3\times3$ matrix: every constant field is divergence free, so the range
condition is vacuous there. Universal cancellation forces Hermiticity on every
nonzero fiber, and cannot force it at zero.

Corollary~B also admits a direct and sharper proof inside the solenoidal
class, using a seven-mode seed and one unequal-length triad per new mode. That
route is retained in Sections~\ref{sec:seeds} and \ref{sec:propagation}
because it yields the optimal $8M-2$ certificate and the projected
conditioning law, neither of which follows from the algebraic reduction above.

\subsection{Pullback to strain--vorticity cancellations}
\label{subsec:corC}

Let $u$ be admissible and write $\omega=\curl u$, $S=\gsym u$. For a
translation-invariant vector multiplier $T(k):\kperp\to\C^{3}$ define
\begin{equation}
\label{eq:Dform}
\Dform_{T}(u)
=\langle\gsym(Tu),\omega\otimes\omega\rangle .
\end{equation}

\begin{corC}[complete strain--vorticity classification]
\label{cor:C}
$\Dform_{T}(u)=0$ for every admissible $u$ if and only if
\begin{equation}
\label{eq:Tfamily}
T=A\curl+B(-\Delta)+\curl B\curl,
\qquad A=A^{T},\quad B=B^{T} ,
\end{equation}
for a unique pair of constant real symmetric matrices. If $T$ is required to
preserve divergence-free output, this reduces to
\begin{equation}
\label{eq:Tsol}
T=c\curl+d(-\Delta),
\qquad c,d\in\R .
\end{equation}
\end{corC}

The proof is in Section~\ref{subsec:pullback}. As in Theorem~A$^{0}$, the
mean-zero hypothesis on $u$ may be dropped and $T(0)$ included as an arbitrary
real $3\times3$ block: $\omega=\curl u$ is always mean zero, and
$\gsym(T(0)\widehat u(0))=0$. Thus \eqref{eq:Tfamily}--\eqref{eq:Tsol}
classify every nonzero block and the zero block is exactly invisible.

\subsection{Sharp finite certificates}
\label{subsec:thmD}

Let $\Lambda\subset\Z^{3}\setminus\{0\}$ be a finite symmetric set of nonzero
modes, write $[k]=\{k,-k\}$, and let $M=|\Lambda/\{\pm1\}|$ be its number of
unoriented modes. The signed Boolean seed $\mathcal B_{7}\cup(-\mathcal
B_{7})$ and the coordinate unit cube $\Lambda_{1}$ are defined in
\eqref{eq:boolean} and \eqref{eq:unitcube}.

\begin{definition}[unequal-triad ordering]
\label{def:unequal}
A finite symmetric set $\Lambda$ has an \emph{unequal-triad ordering} if it
contains the signed Boolean seed and its remaining unoriented classes can be
ordered as $[t_{1}],\dots,[t_{M-7}]$ so that at step $j$ there are signed modes
$q_{j},r_{j}$ in the seed or in earlier classes with
\[
t_{j}=q_{j}+r_{j},
\qquad
q_{j}\times r_{j}\ne0,
\qquad
|q_{j}|\ne|r_{j}| .
\]
The full six-mode support $\{\pm t_{j},\pm q_{j},\pm r_{j}\}$ of each
witnessing triad is required to lie in $\Lambda$.
\end{definition}

\begin{definition}[two-triad ordering]
\label{def:twotriad}
A finite symmetric set $\Lambda$ has a \emph{two-triad ordering} if it contains
$\Lambda_{1}$ and its remaining unoriented classes can be ordered as
$[t_{1}],\dots,[t_{M-13}]$ so that at step $j$ there are earlier signed modes
$q_{j},r_{j},q'_{j},r'_{j}$ with
\[
t_{j}=q_{j}+r_{j}=q'_{j}+r'_{j},
\qquad
q_{j}\times r_{j}\ne0,
\qquad
q'_{j}\times r'_{j}\ne0,
\qquad
r_{j}-q_{j}\not\parallel r'_{j}-q'_{j} .
\]
Again the full six-mode support of both witnessing triads must lie in
$\Lambda$.
\end{definition}

\begin{thmD}[sharp finite certificates]
\label{thm:D}
Let $\Lambda$ be a finite symmetric set with $M$ unoriented modes.
\begin{enumerate}
\item If $\Lambda$ has an unequal-triad ordering, then a general complex
solenoidal symbol on $\Lambda$ has $8M$ real parameters, and there exists a
family of exactly $8M-2$ real triad-supported field evaluations certifying the
energy--helicity form \eqref{eq:scalarform}.
\item Under the same ordering hypothesis, there exists a certifying family of
exactly $4M-2$ evaluations in the Hermitian solenoidal subclass, and this
cardinality is minimal.
\item If $\Lambda$ has a two-triad ordering, then a general
unrestricted-output symbol on $\Lambda$ has $12M$ real parameters, and there
exists a family of exactly $12M-12$ evaluations certifying
\eqref{eq:generator}.
\end{enumerate}
No scalar test family with fewer evaluations can work in any of these
settings, because one
real field evaluation supplies at most one real linear equation while the
unavoidable kernels have dimensions $2$, $2$ and $12$ respectively.
\end{thmD}

The theorem makes no claim for a finite symmetric set lacking one of the two
orderings. Its proof is in Section~\ref{subsec:thmDproof}. For every
coordinate box $\Lambda_{N}=\{k\ne0:|k_{j}|\le N\}$ with $N\ge1$ both ordering
hypotheses hold, and
\begin{equation}
\label{eq:boxcount}
M_{N}=4N^{3}+6N^{2}+3N ,
\end{equation}
so that the sharp counts are
\begin{equation}
\label{eq:boxtests}
\begin{array}{c|c}
\text{symbol class}&\text{number of tests}\\ \hline
\text{Hermitian solenoidal}&16N^{3}+24N^{2}+12N-2\\
\text{complex solenoidal}&32N^{3}+48N^{2}+24N-2\\
\text{unrestricted output}&48N^{3}+72N^{2}+36N-12 .
\end{array}
\end{equation}
These families have the optimal $\Theta(N^{3})$ size rather than the
$\Theta(N^{6})$ cost of scanning a dense collection of triads. They classify
the nonzero blocks; if a zero-mode block is appended, no field evaluation can
constrain its nine real parameters, so sharpness is understood modulo that
exactly invisible block.

\section{Sufficiency: why the displayed generators cancel}
\label{sec:sufficiency}

The sufficiency half of Theorem~A is analytic and independent of every finite
certificate. Let $w$ be admissible and $A=A^{T}$ constant. Then
\begin{equation}
\label{eq:suff-A}
\Qform_{A}(w)
=\int_{\T^{3}}(Aw)\cdot(w\cdot\nabla)w\,dx
=\frac12\int_{\T^{3}}w\cdot\nabla(w\cdot Aw)\,dx=0 ,
\end{equation}
the last step because $\dv w=0$.

For the first-order family put $\omega=\curl w$ and $v=(w\cdot\nabla)w$.
Since $\dv w=\dv\omega=0$,
\begin{equation}
\label{eq:curlconv}
\curl v=(w\cdot\nabla)\omega-(\omega\cdot\nabla)w .
\end{equation}
Using self-adjointness of $\curl$ and \eqref{eq:curlconv}, for constant
$B=B^{T}$,
\begin{equation}
\label{eq:suff-B}
\begin{aligned}
\Qform_{B\curl+\curl B}(w)
&=\langle B\omega,v\rangle+\langle Bw,\curl v\rangle\\
&=\langle B\omega,(w\cdot\nabla)w\rangle
 +\langle Bw,(w\cdot\nabla)\omega\rangle
 -\langle Bw,(\omega\cdot\nabla)w\rangle .
\end{aligned}
\end{equation}
By symmetry of $B$ the first two terms are together
$\int(w\cdot\nabla)(w\cdot B\omega)\,dx$, which vanishes because $\dv w=0$.
The third is
\begin{equation}
\label{eq:suff-B2}
-\frac12\int_{\T^{3}}\omega\cdot\nabla(w\cdot Bw)\,dx=0 ,
\end{equation}
because $\dv\omega=0$. Hence $\Qform_{A+B\curl+\curl B}\equiv0$ for all
constant symmetric $A,B$, which is the sufficiency half of Theorem~A, and by
\eqref{eq:zeromean}--\eqref{eq:commute} also of Theorem~A$^{0}$.

When $A=cI$, \eqref{eq:suff-A} is the kinetic-energy cancellation; when
$B=bI$, \eqref{eq:suff-B} is the helicity cancellation, with the normalization
$B\curl+\curl B=2b\curl$.

\section{Triad reduction}
\label{sec:triad}

\subsection{Isolating a complex triad coefficient}
\label{subsec:phase}

Real test fields carry a reality constraint, so it is not immediate that they
see the full complex triad coefficient of a possibly non-Hermitian symbol.
They do, and the mechanism is a two-element phase family.

\begin{lemma}[polarization by the phases $1$ and $i$]
\label{lem:phase}
Let $p+q+r=0$ be a non-collinear lattice triad and choose real transverse
polarizations $A_{0}\perp p$, $U\perp q$, $V\perp r$. Define the complex triad
coefficient
\begin{equation}
\label{eq:triadcoeff}
\begin{aligned}
\mathfrak T_{R}(p,q,r;A_{0},U,V)=i\bigl[{}
&R(p)A_{0}\cdot\bigl((U\cdot r)V+(V\cdot q)U\bigr)\\
{}+{}&R(q)U\cdot\bigl((V\cdot p)A_{0}+(A_{0}\cdot r)V\bigr)\\
{}+{}&R(r)V\cdot\bigl((A_{0}\cdot q)U+(U\cdot p)A_{0}\bigr)
\bigr].
\end{aligned}
\end{equation}
For phases $\alpha,\beta,\gamma\in\C$ let $w$ be the real six-mode field
determined by
$\widehat w(p)=\alpha A_{0}$, $\widehat w(q)=\beta U$,
$\widehat w(r)=\gamma V$ and $\widehat w(-k)=\overline{\widehat w(k)}$. Then
\begin{equation}
\label{eq:Qphase}
\Qform_{R}(w)
=2\operatorname{Re}\bigl(\alpha\beta\gamma\,
\mathfrak T_{R}(p,q,r;A_{0},U,V)\bigr).
\end{equation}
Consequently, if $\Qform_{R}$ vanishes for the two phase products
$\alpha\beta\gamma=1$ and $\alpha\beta\gamma=i$, then
\[
\mathfrak T_{R}(p,q,r;A_{0},U,V)=0 .
\]
\end{lemma}

\begin{proof}
Expanding the Fourier coefficient of $(w\cdot\nabla)w$ and grouping the terms
by recipient mode gives the three summands in \eqref{eq:triadcoeff}. It remains
to check that no other signed zero-sum triple contributes.

Non-collinearity means $p\times q\ne0$, so $p,q,r$ are nonzero and pairwise
non-collinear: if for instance $p\parallel r$, then $q=-p-r\parallel p$, a
contradiction. Consider any three modes from $\{\pm p,\pm q,\pm r\}$ summing to
zero. If an unoriented mode is repeated, then either an opposite pair cancels
and leaves a nonzero third mode, or the relation has the form $\pm2x\pm y=0$
or $\pm3x=0$; all alternatives are impossible by nonvanishing and pairwise
non-collinearity. So the three unoriented modes are $p,q,r$. Writing their
signs as $\varepsilon,\eta,\theta\in\{\pm1\}$ and using $r=-p-q$ gives
\[
(\varepsilon-\theta)p+(\eta-\theta)q=0 ,
\]
and linear independence of $p,q$ forces $\varepsilon=\eta=\theta$. Hence the
only signed zero-sum triples are permutations of $(p,q,r)$ and $(-p,-q,-r)$.
The positive triad contributes $\alpha\beta\gamma\mathfrak T_{R}$ and, by
reality compatibility, the negative one contributes its complex conjugate;
their sum is \eqref{eq:Qphase}.

Writing $\mathfrak T_{R}=x+iy$, the phase product $1$ gives $2x=0$ and the
phase product $i$ gives $-2y=0$.
\end{proof}

For a fixed polarization triple the eight choices
$(\alpha,\beta,\gamma)\in\{1,i\}^{3}$ used by the certificate programs have
total products in $\{1,i,-1,-i\}$; they are redundant up to row sign and span
exactly the same two real constraints as the products $1$ and $i$.

Lemma~\ref{lem:phase} is the reason a non-Hermitian symbol is detectable at
all: the skew-adjoint sector contributes to \eqref{eq:intro-Q}, even though it
would be invisible in a real quadratic form.

\subsection{The raw recipient map}
\label{subsec:raw}

The convective output received at the mode $p$, up to a fixed nonzero scalar,
is
\begin{equation}
\label{eq:rawN}
N_{q,r}(U,V)=(U\cdot r)V+(V\cdot q)U .
\end{equation}

\begin{lemma}[raw image line and two-triad determination]
\label{lem:rawline}
Let $p+q+r=0$ be non-collinear and suppose $D$ satisfies the universal
cancellation \eqref{eq:universal} and vanishes at $q$ and at $r$. Then
\begin{equation}
\label{eq:imageline}
D(p)\bigl(p_{\C}^{\perp}\bigr)\subset\spn_{\C}\{r-q\} .
\end{equation}
One raw triad therefore has real rank eight on the twelve-parameter target
block. If a second known-parent triad at the same recipient has difference
direction not parallel to $r-q$, the two triads force $D(p)=0$ and their
stacked target block has real rank twelve.
\end{lemma}

\begin{proof}
Put $z=q\times r$, $Z=|z|$ and $Jx=z\times x$, and write
$U=u_{0}z+\alpha Jq$, $V=v_{0}z+\beta Jr$. A direct computation gives
\begin{equation}
\label{eq:Nexplicit}
N_{q,r}(U,V)
=Z^{2}\bigl[(\alpha v_{0}-\beta u_{0})z+\alpha\beta J(r-q)\bigr],
\end{equation}
so that
\begin{equation}
\label{eq:Nrange}
\range N_{q,r}=\spn\{z,J(r-q)\} .
\end{equation}
Since $z\times J(r-q)=-Z^{2}(r-q)$, the orthogonal complement of this range is
exactly $\spn\{r-q\}$.

Lemma~\ref{lem:phase}, with the two parent terms in \eqref{eq:triadcoeff}
equal to zero, gives
\begin{equation}
\label{eq:rawconstraint}
D(p)A_{0}\cdot N_{q,r}(U,V)=0
\end{equation}
for all polarizations, which is \eqref{eq:imageline}.

A general complex $3\times2$ target block has twelve real parameters. One raw
triad imposes rank eight and leaves the four-real-dimensional family
\begin{equation}
\label{eq:residual}
D(p)=(r-q)\otimes\ell,
\qquad
\ell:p_{\C}^{\perp}\to\C .
\end{equation}
If a second known-parent decomposition of the same target produces an image
line not parallel to $r-q$, the two conditions \eqref{eq:imageline} intersect
only at zero, so $D(p)=0$ and the stacked block has rank twelve.
\end{proof}

\subsection{Leray projection and the isosceles defect}
\label{subsec:projected}

In the solenoidal class the recipient is projected:
\begin{equation}
\label{eq:projN}
N_{p}(U,V)=P_{p^{\perp}}\bigl[(U\cdot r)V+(V\cdot q)U\bigr].
\end{equation}

\begin{lemma}[projected one-triad determination]
\label{lem:projected}
Suppose $D$ is solenoidal, satisfies \eqref{eq:universal}, and vanishes at $q$
and at $r$. If $|q|\ne|r|$, then $D(p)=0$. If $|q|=|r|$, the target ranks are
exactly four for an arbitrary complex solenoidal block and three for a
Hermitian block.
\end{lemma}

\begin{proof}
With the notation of Lemma~\ref{lem:rawline}, the choices $(U,V)=(z,Jr)$ and
$(U,V)=(Jq,Jr)$ give
\begin{equation}
\label{eq:twooutputs}
N_{p}=-|z|^{2}z,
\qquad
N_{p}=|z|^{2}P_{p^{\perp}}J(r-q).
\end{equation}
The second vector has a nonzero in-plane component precisely when
\begin{equation}
\label{eq:inplane}
J(r-q)\cdot Jp=|z|^{2}\bigl(|q|^{2}-|r|^{2}\bigr)\ne0 .
\end{equation}
So if $|q|\ne|r|$ the two outputs span $p^{\perp}$ and determine all eight
real parameters of an arbitrary complex $2\times2$ target symbol. If
$|q|=|r|$, the projected range is only $\spn\{z\}$: one complex output row is
annihilated and the other remains free.

Concretely, in an orthonormal output basis whose first vector is parallel to
$z$, the unequal-length case tests both complex output rows of $D(p)$, while
the equal-length case tests only the first. The latter gives two complex,
hence four real, constraints on a general block. For a Hermitian block the
first row consists of one real diagonal entry and one arbitrary complex
off-diagonal entry, hence three real constraints.
\end{proof}

The exact local ranks are therefore
\begin{equation}
\label{eq:localranks}
\begin{array}{c|c|c}
&|q|\ne|r|&|q|=|r|\\ \hline
\text{arbitrary complex solenoidal block}&8&4\\
\text{Hermitian solenoidal block}&4&3 .
\end{array}
\end{equation}
The raw map \eqref{eq:rawN} has no isosceles degeneracy at all. Equal parent
lengths become exceptional only after the longitudinal direction has been
discarded by the Leray projection. Section~\ref{sec:singular} computes the
singular value that vanishes in the limit and shows that it is controlled by
the logarithmic separation of the two parent shells.

\section{Exact finite seeds}
\label{sec:seeds}

The necessity proofs begin with exact finite-dimensional certificates. Each
script symbol block is expressed in the integer, non-normalized frame covered
by Lemma~\ref{lem:coordinates}. Test fields use two integer transverse basis
polarizations at each mode, the phase choices described after
Lemma~\ref{lem:phase}, and conjugate negative-mode amplitudes. Expanding
\eqref{eq:intro-Q} produces a real rational matrix in the symbol parameters,
and Lemma~\ref{lem:coordinates} transfers its rank and nullity to the
normalized convention used here.

\subsection{The solenoidal seed}
\label{subsec:solseed}

Use the seven unoriented Boolean modes
\begin{equation}
\label{eq:boolean}
\mathcal B_{7}=
\{e_{1},e_{2},e_{3},e_{1}+e_{2},e_{1}+e_{3},e_{2}+e_{3},e_{1}+e_{2}+e_{3}\}.
\end{equation}
A general complex $2\times2$ block has eight real parameters, so the seed
symbol space has dimension $56$. Write $\mathsf M_{\mathcal B_{7}}$ for the
resulting exact matrix. It has
\begin{equation}
\label{eq:solseedrank}
720\text{ nonzero rows},
\qquad
\rank=54,
\qquad
\nullity=2 .
\end{equation}
A fixed $54\times54$ minor has determinant residues
\begin{equation}
\label{eq:solminor}
190{,}194
\quad\text{and}\quad
891{,}695{,}525
\end{equation}
modulo $1{,}000{,}003$ and $1{,}000{,}000{,}007$ respectively. All row
denominators lie in $\{1,2,3,6\}$, so neither prime annihilates a denominator,
and either nonzero residue proves $\rank\ge54$ over $\Q$. The two linearly
independent kernel vectors given by the restrictions of $I$ and $\Ck$ prove the
matching upper bound. Rank--nullity then yields the exact kernel
\emph{equality}
\begin{equation}
\label{eq:seed-kernel-sol}
\ker\mathsf M_{\mathcal B_{7}}
=\Bigl\{
\bigl(\left.(cI+d\Ck)\right|_{\kperp}\bigr)_{k\in\mathcal B_{7}\cup(-\mathcal B_{7})}
:c,d\in\R
\Bigr\}.
\end{equation}
Consequently every reality-compatible solenoidal symbol annihilating the seed
tests agrees on the signed seed with one and only one symbol $cI+d\Ck$.

In the Hermitian subclass the same seed gives
\begin{equation}
\label{eq:hermseed}
592\text{ nonzero rows},
\qquad
28\text{ parameters},
\qquad
\rank=26 .
\end{equation}
Since $I$ and $\Ck$ are Hermitian, rank $26$ and nullity $2$ give the same
kernel \eqref{eq:seed-kernel-sol} restricted to the Hermitian parameter space,
so the Hermitian seed determines the same unique pair $c,d$.

An exhaustive exact search through all $4{,}095$ subsets of one through six
modes of the coordinate unit cube shows that none attains the required
Hermitian rank $4M-2$. Exactly four seven-mode subsets attain rank $26$: the
nonzero subset sums of $e_{1},\pm e_{2},\pm e_{3}$, modulo total sign. This is
minimality within a finite domain only, and is not a statement about arbitrary
six-mode configurations elsewhere in the lattice.

\subsection{The unrestricted-output seed}
\label{subsec:rawseed}

Use the full coordinate unit cube
\begin{equation}
\label{eq:unitcube}
\Lambda_{1}=\{k\in\Z^{3}\setminus\{0\}:|k_{j}|\le1\},
\end{equation}
which contains thirteen modes modulo total sign. A general complex $3\times2$
block has twelve real parameters, so the parameter space has dimension $156$.
Enumerating the $44$ non-collinear triads up to ordering and total sign gives an
integer matrix
\begin{equation}
\label{eq:M1}
\mathsf M_{1}\in\Z^{680\times156},
\qquad
\rank\mathsf M_{1}=144,
\qquad
\nullity\mathsf M_{1}=12 .
\end{equation}
The twelve explicit generators obtained from bases of $\Sym(3,\R)$ for $A$ and
$B$ are linearly independent and annihilate every row, so rank--nullity again
gives an exact kernel equality,
\begin{equation}
\label{eq:seed-kernel-raw}
\ker\mathsf M_{1}
=\Bigl\{
\bigl(\left.(A+B\Ck+\Ck B)\right|_{\kperp}\bigr)_{k\in\Lambda_{1}}
:A,B\in\Sym(3,\R)
\Bigr\}.
\end{equation}
Proposition~\ref{prop:uniqueness} uses only the signed coordinate-axis modes,
which belong to $\Lambda_{1}$, so the matrices $A,B$ in
\eqref{eq:seed-kernel-raw} are unique. Thus every unrestricted-output symbol
annihilating the unit-cube tests agrees there with one and only one generator
of the form \eqref{eq:generator}.

A fixed $144\times144$ minor has residues
\begin{equation}
\label{eq:rawminor}
272{,}144
\quad\text{and}\quad
290{,}707{,}620
\end{equation}
modulo the same two primes. An independent convolution-first implementation
reconstructs the $680$ rows and fixes the canonical digest
\begin{center}
\texttt{a4fc3b45a6fc94defd38d992dbbe27f4413254b3af307e41d1c7e25839e9c30d}
\end{center}
(SHA-256 of the canonically ordered row data).

The seven-mode Boolean seed does \emph{not} suffice in this larger class: it
has $84$ parameters, rank $66$ and nullity $18$. The six additional
finite-network freedoms disappear on the full unit cube. This is a concrete
illustration of the difference between invariants of a restricted triad
network and universal ones.

\subsection{Certificate summary}
\label{subsec:certsummary}

\begin{center}
\begin{tabular}{lrrrrr}
\hline
Class & Seed & Parameters & Nonzero rows & Rank & Kernel\\
\hline
Hermitian, solenoidal output & $7$ modes & $28$ & $592$ & $26$ & $2$\\
Complex, solenoidal output & $7$ modes & $56$ & $720$ & $54$ & $2$\\
Complex, unrestricted output & $13$ modes & $156$ & $680$ & $144$ & $12$\\
\hline
\end{tabular}
\end{center}

The construction above records the mathematical generation of these matrices
and the modular-minor argument used in the proof. The programs are a
reproducible certificate; the finite objects, and the implication from their
ranks to the theorems, are stated here explicitly and do not depend on reading
the code.

\section{Propagation and proofs of the classification theorems}
\label{sec:propagation}

Throughout this section we use the target/parent/recipient convention
\eqref{eq:targetconv}: the new mode is $t=q+r$, the recipient at which the
local lemmas act is $p=-t$, and reality transports $D(-t)=0$ to $D(t)=0$.

\subsection{Solenoidal one-triad propagation}
\label{subsec:solprop}

By the kernel identity \eqref{eq:seed-kernel-sol} there is a unique $cI+d\Ck$
agreeing with $R$ on $\mathcal B_{7}\cup(-\mathcal B_{7})$. Subtract it and
call the difference $D$, so that $D$ is a reality-compatible solenoidal symbol
satisfying \eqref{eq:universal} and vanishing on the fourteen signed seed
modes. We show $D\equiv0$.

Six unequal-length steps recover the missing unit-cube modes. First
\begin{equation}
\label{eq:boot1}
\begin{aligned}
e_{1}+e_{2}-e_{3}&=(e_{1}+e_{2})-e_{3},\\
e_{1}-e_{2}+e_{3}&=(e_{1}+e_{3})-e_{2},\\
e_{1}-e_{2}-e_{3}&=e_{1}-(e_{2}+e_{3}),
\end{aligned}
\end{equation}
whose parents have squared lengths $2$ and $1$; then
\begin{equation}
\label{eq:boot2}
\begin{aligned}
e_{1}-e_{2}&=(e_{1}-e_{2}+e_{3})-e_{3},\\
e_{1}-e_{3}&=(e_{1}-e_{2}-e_{3})+e_{2},\\
e_{2}-e_{3}&=(e_{1}+e_{2}-e_{3})-e_{1},
\end{aligned}
\end{equation}
whose parents have squared lengths $3$ and $1$. In each equality the left-hand
side is the target $t$ and the two right-hand vectors are $q,r$; the parents
are non-collinear with distinct lengths, so Lemma~\ref{lem:projected} gives
$D(-t)=0$, and reality gives $D(t)=0$. Permuting coordinates covers the
remaining unit-cube classes.

All axis modes follow from
\begin{equation}
\label{eq:axis1}
2e_{i}+e_{j}=e_{i}+(e_{i}+e_{j}),
\qquad
2e_{i}=(2e_{i}+e_{j})-e_{j},
\end{equation}
and inductively, for $n\ge2$,
\begin{equation}
\label{eq:axis2}
(n+1)e_{i}+e_{j}=ne_{i}+(e_{i}+e_{j}),
\qquad
(n+1)e_{i}=\bigl((n+1)e_{i}+e_{j}\bigr)-e_{j}.
\end{equation}
Again each left-hand side is a target and each right-hand side its known-parent
decomposition, with non-collinear parents of distinct lengths.

For a remaining non-axis target $t$ outside the recovered coordinate unit
cube, so that $|t|_{1}\ge3$, induct on $|t|_{1}$. Choose a nonzero coordinate,
set $s=\sgn(t_{j})e_{j}$, and write
\begin{equation}
\label{eq:generalstep}
t=(t-s)+s .
\end{equation}
For example, $t=(2,1,0)$ is obtained from $(1,1,0)+e_{1}$; the same
construction applies to every non-axis target. Indeed, $t-s$ has
$\ell^{1}$-norm $|t|_{1}-1\ge2$, so its squared Euclidean norm is at least
$2$, while $|s|^{2}=1$. The two parents are non-collinear: $t$ has a nonzero
coordinate other than the one used to define $s$, and that coordinate remains
in $t-s$. Thus the parents have distinct lengths and Lemma~\ref{lem:projected}
at $p=-t$ determines $D(-t)$; reality determines $D(t)$. Since the parents
lie in earlier $\ell^{1}$-shells or in the recovered seed, this is a genuine
induction. It reaches every non-axis mode, and together with the axis steps
above proves $D\equiv0$ and $R=cI+d\Ck$. This proves Corollary~B directly,
without passing through Theorem~A.

\subsection{Raw two-triad propagation, and the proof of Theorem~A}
\label{subsec:rawprop}

For Theorem~A the kernel identity \eqref{eq:seed-kernel-raw} supplies the
unique generator \eqref{eq:generator} agreeing with $R$ on the unit cube.
Subtract it and again call the difference $D$, a reality-compatible symbol with
unrestricted output that satisfies \eqref{eq:universal} and vanishes on
$\Lambda_{1}\cup(-\Lambda_{1})$.

Induct on $|t|_{1}$, starting at shell $3$: every non-axis mode in shells $1$
and $2$ already lies in the coordinate unit cube. Within each shell, process
all non-axis targets before the axis targets.

Let $t$ be a non-axis target with $|t|_{1}=n\ge3$, choose two nonzero
coordinates $i,j$, and set $s_{i}=\sgn(t_{i})e_{i}$, $s_{j}=\sgn(t_{j})e_{j}$.
Use the two parent decompositions
\begin{equation}
\label{eq:twoparents}
t=(t-s_{i})+s_{i}=(t-s_{j})+s_{j}.
\end{equation}
Here $t-s_{i}$ and $t-s_{j}$ lie in shell $n-1$ while $s_{i},s_{j}$ lie in the
unit cube, so all four parents are already known; each pair is non-collinear
because $t$ has another nonzero coordinate. At the common recipient $p=-t$,
the two surviving image lines from \eqref{eq:imageline} have directions
\begin{equation}
\label{eq:difdirs}
d_{i}=2s_{i}-t,
\qquad
d_{j}=2s_{j}-t .
\end{equation}
They are non-parallel. Indeed, if $t$ has a third nonzero coordinate,
comparison in that coordinate forces the proportionality factor to be one,
whereas the $i$ coordinate would then require
$2\sgn(t_{i})-t_{i}=-t_{i}$, which is impossible. If $t$ has exactly two
nonzero coordinates, of magnitudes $a,b$, then
\[
d_{i}\times d_{j}=2(2-a-b)\,s_{i}\times s_{j}\ne0
\]
because $a+b=|t|_{1}\ge3$. By Lemma~\ref{lem:rawline} the two raw triads force
$D(-t)=0$, and reality gives $D(t)=0$.

Once the non-axis part of shell $n$ is known, take the positive axis target
$t=ne_{i}$ with $n\ge2$ and use
\begin{equation}
\label{eq:axisraw}
\begin{aligned}
ne_{i}&=\bigl((n-1)e_{i}+e_{j}\bigr)+(e_{i}-e_{j}),\\
ne_{i}&=\bigl((n-1)e_{i}+e_{\ell}\bigr)+(e_{i}-e_{\ell}),
\end{aligned}
\end{equation}
with $j\ne\ell$ both different from $i$. Here $(n-1)e_{i}+e_{j}$ and
$(n-1)e_{i}+e_{\ell}$ are non-axis modes of the current shell, already handled,
while $e_{i}-e_{j}$ and $e_{i}-e_{\ell}$ belong to the unit cube. Hence every
parent in \eqref{eq:axisraw} is known at this stage, including when $n=2$. The
difference directions
\[
(2-n)e_{i}-2e_{j},
\qquad
(2-n)e_{i}-2e_{\ell}
\]
are non-parallel, so Lemma~\ref{lem:rawline} at $p=-ne_{i}$ gives
$D(-ne_{i})=0$ and reality gives $D(ne_{i})=0$; the negative axis mode follows
simultaneously. This completes the induction, so $D\equiv0$ and
\eqref{eq:generator} holds. Together with Section~\ref{sec:sufficiency} and
Proposition~\ref{prop:uniqueness}, Theorem~A is proved, and with
\eqref{eq:zeromean}--\eqref{eq:commute} so is Theorem~A$^{0}$.

\subsection{Coverage check for the propagation routes}
\label{subsec:propcheck}

The case distinctions above can be checked independently of the rank
certificates. For the solenoidal route, the bootstrap identities recover the
three missing body diagonals and then the three missing face diagonals of the
unit cube. The axis identities recover each positive axis in increasing
$\ell^{1}$-order, and the general step recovers every target with at least two
nonzero coordinates. Its parents have $\ell^{1}$-norm one and one less than
that of the target, so the induction is well founded. Signs are transported by
reality. For the raw route, every non-axis target has at least two nonzero
coordinates and therefore admits the two decompositions in
\eqref{eq:twoparents}; after these targets are processed, the decompositions in
\eqref{eq:axisraw} recover each axis. Thus no target is omitted by either
ordering.

As an exact finite audit of these assertions, the verification programs replay
both orderings in the coordinate box $|k_j|\le10$. The solenoidal replay
reaches all $9{,}260$ signed nonzero modes. The raw replay checks all $4{,}617$
target blocks, verifies that both parent pairs are known when used, and checks
that the two difference directions have nonzero cross product and that the
stacked target block has exact rank twelve in every case. These computations
are finite checks of the stated formulas; the preceding inductions extend the
argument from the checked box to all of $\Z^{3}\setminus\{0\}$.

\subsection{Structural comparison}
\label{subsec:comparison}

The two propagation arguments are summarized by one geometric dichotomy:
\begin{equation}
\label{eq:dichotomy}
\begin{array}{c|c|c}
\text{output}&\text{one triad leaves}&\text{what kills the target}\\ \hline
\C^{3}&\text{one complex image line}&\text{a second non-parallel line}\\
p^{\perp}&\text{no freedom if }|q|\ne|r|&\text{one unequal-length triad}.
\end{array}
\end{equation}
Projection makes generic propagation stronger --- one triad suffices instead of
two --- but it creates the exact isosceles defect, which is the subject of
Section~\ref{sec:singular}.

\subsection{Proof of Theorem~D}
\label{subsec:thmDproof}

Definitions~\ref{def:unequal} and \ref{def:twotriad} isolate exactly what the
preceding inductions use.

For an unequal-triad ordering, exact elimination selects $54$ independent rows
on the complex solenoidal seed. At each later target, Lemma~\ref{lem:projected}
supplies eight rows whose restriction to the new target block has rank eight.
Ordering rows and columns by propagation makes the matrix block lower
triangular, of rank
\begin{equation}
\label{eq:count8}
54+8(M-7)=8M-2 .
\end{equation}
Restricting the same construction to Hermitian blocks selects $26$ seed rows
and four rows at each later target, giving rank
\begin{equation}
\label{eq:count4}
26+4(M-7)=4M-2 ,
\end{equation}
which proves clause 2 as the Hermitian specialization of the same triangular
construction.

For a two-triad ordering, select the $144$ independent unit-cube rows. The two
non-parallel image lines at each subsequent target supply a rank-twelve
diagonal block, so the total rank is
\begin{equation}
\label{eq:count12}
144+12(M-13)=12M-12 .
\end{equation}

All selected fields are supported in $\Lambda$ by the corresponding ordering
definition. The energy--helicity plane lies in the kernel of the first two
matrices and the twelve symbols \eqref{eq:generator} lie in the kernel of the
third, so the displayed ranks give exactly those kernels. Conversely, one real
field evaluation supplies at most one real linear equation, so the parameter
dimensions $8M$, $4M$, $12M$ and the kernel dimensions $2$, $2$, $12$ give the
matching lower bounds. This proves Theorem~D.

For a coordinate box $\Lambda_{N}$, the bootstrap steps
\eqref{eq:boot1}--\eqref{eq:boot2}, the axis steps
\eqref{eq:axis1}--\eqref{eq:axis2} and the induction \eqref{eq:generalstep} all
remain inside $\Lambda_{N}$ and give an unequal-triad ordering; the two-parent
construction \eqref{eq:twoparents}--\eqref{eq:axisraw} likewise remains inside
$\Lambda_{N}$ and gives a two-triad ordering. Counting unoriented modes gives
\eqref{eq:boxcount} and hence \eqref{eq:boxtests}. Finite symmetric sets
without one of these orderings lie outside the theorem.

\section{Singular values and quantitative rigidity}
\label{sec:singular}

Set
\[
\lambda=|q|^{2},
\quad
\mu=|r|^{2},
\quad
\kappa=|p|^{2},
\quad
z=q\times r,
\quad
Z=|z|,
\quad
Jv=z\times v .
\]
All singular values in Sections~\ref{subsec:rawspec}--\ref{subsec:routes} are
\emph{local}: they refer to the Euclidean structures induced by orthonormal
frames and unit tensor-product polarizations, for the maps
\begin{equation}
\label{eq:localmaps}
\mathsf L^{\mathrm{raw}}_{p;q,r}(U\otimes V)
=(U\cdot r)V+(V\cdot q)U,
\qquad
\mathsf L^{\mathrm{sol}}_{p;q,r}
=P_{p^{\perp}}\mathsf L^{\mathrm{raw}}_{p;q,r} .
\end{equation}
Their singular spectra are invariant under orthonormal changes of basis.
Lemma~\ref{lem:coordinates} explains why the non-normalized certificate
coordinates preserve the corresponding ranks but not these numerical values.

\subsection{The raw interaction spectrum}
\label{subsec:rawspec}

\begin{proposition}[raw singular spectrum]
\label{prop:raw}
For the unprojected map \eqref{eq:rawN}, the two nonzero squared singular
values are
\begin{equation}
\label{eq:rawspec}
\widetilde\sigma_{\mathrm n}^{2}
=\frac{Z^{2}(\lambda+\mu)}{\lambda\mu},
\qquad
\widetilde\sigma_{\mathrm d}^{2}
=\frac{Z^{2}|r-q|^{2}}{\lambda\mu} .
\end{equation}
Both are positive for every non-collinear triad.
\end{proposition}

\begin{proof}
The vectors
\[
n=\frac zZ,
\qquad
t_{q}=\frac{Jq}{Z\sqrt\lambda},
\qquad
t_{r}=\frac{Jr}{Z\sqrt\mu},
\qquad
t_{d}=\frac{J(r-q)}{Z|r-q|}
\]
are orthonormal frames in the two parent planes and in the raw
two-dimensional image. Formula \eqref{eq:Nexplicit}, applied first to the
unnormalized vectors $z,Jq,Jr$, gives
\[
N_{q,r}(z,z)=0,
\quad
N_{q,r}(z,Jr)=-Z^{2}z,
\quad
N_{q,r}(Jq,z)=Z^{2}z,
\quad
N_{q,r}(Jq,Jr)=Z^{2}J(r-q).
\]
Hence in the ordered tensor-product input basis
$(n\otimes n,\;n\otimes t_{r},\;t_{q}\otimes n,\;t_{q}\otimes t_{r})$ and the
output basis $(n,t_{d})$, the raw interaction has matrix
\begin{equation}
\label{eq:rawmatrix}
\begin{pmatrix}
0&-Z/\sqrt\mu&Z/\sqrt\lambda&0\\
0&0&0&Z|r-q|/(\sqrt\lambda\sqrt\mu)
\end{pmatrix}.
\end{equation}
Its two rows are orthogonal, and their squared Euclidean norms are exactly the
quantities in \eqref{eq:rawspec}.
\end{proof}

\subsection{The projected interaction spectrum}
\label{subsec:projspec}

\begin{proposition}[projected spectrum and target observation]
\label{prop:projected}
For the Leray-projected map \eqref{eq:projN}, the squared singular values are
\begin{equation}
\label{eq:projspec}
\sigma_{\mathrm n}^{2}
=\frac{Z^{2}(\lambda+\mu)}{\lambda\mu},
\qquad
\sigma_{\mathrm t}^{2}
=\frac{Z^{2}(\lambda-\mu)^{2}}{\kappa\lambda\mu} .
\end{equation}
The projection leaves the normal channel unchanged and multiplies the second
raw singular value by
\begin{equation}
\label{eq:defectfactor}
\frac{|\lambda-\mu|}{\sqrt\kappa\,|r-q|} .
\end{equation}
Let $\mathsf O_{p;q,r}(D(p))$ denote the complete local target observation
formed from the coefficients $D(p)A_{0}\cdot N_{p}(U,V)$ over orthonormal input
polarizations. If, in the orthonormal frame $(n,t_{p})$,
\[
D(p)=\begin{pmatrix}a&c\\ \overline c&b\end{pmatrix},
\qquad a,b\in\R,\ c\in\C,
\]
is Hermitian, then
\begin{equation}
\label{eq:obsnorm}
\|\mathsf O_{p;q,r}(D(p))\|_{\mathrm F}^{2}
=\sigma_{\mathrm n}^{2}(a^{2}+|c|^{2})
+\sigma_{\mathrm t}^{2}(b^{2}+|c|^{2}),
\end{equation}
and in particular
\begin{equation}
\label{eq:obsbound}
\|\mathsf O_{p;q,r}(D(p))\|_{\mathrm F}
\ge\min(\sigma_{\mathrm n},\sigma_{\mathrm t})\,\|D(p)\|_{\mathrm{HS}} .
\end{equation}
For an arbitrary complex solenoidal block each singular value in
\eqref{eq:projspec} occurs with real multiplicity four in
$\mathsf O_{p;q,r}$.
\end{proposition}

\begin{proof}
Replace $t_{d}$ in Proposition~\ref{prop:raw} by the normalized tangent vector
$t_{p}=Jp/(Z\sqrt\kappa)$. The identities
\eqref{eq:twooutputs}--\eqref{eq:inplane}, together with
\[
N_{p}(Jq,Jr)=Z^{2}\frac{\lambda-\mu}{\kappa}Jp ,
\]
show that in the same input basis and the output basis $(n,t_{p})$ the
projected interaction has matrix
\begin{equation}
\label{eq:projmatrix}
\begin{pmatrix}
0&-Z/\sqrt\mu&Z/\sqrt\lambda&0\\
0&0&0&Z(\lambda-\mu)/(\sqrt\kappa\sqrt\lambda\sqrt\mu)
\end{pmatrix}.
\end{equation}
Its rows are again orthogonal, and their squared norms give
\eqref{eq:projspec}. Dividing the absolute value of the last entry of
\eqref{eq:projmatrix} by the corresponding entry of \eqref{eq:rawmatrix} gives
\eqref{eq:defectfactor}, which exhibits the isosceles defect explicitly.
Contracting the two orthogonal output rows against the two rows of the
Hermitian matrix $D(p)$ gives \eqref{eq:obsnorm} and hence
\eqref{eq:obsbound}. Without the Hermitian relation each of the two complex
rows of $D(p)$ contains two complex coordinates, and separating real and
imaginary parts gives real multiplicity four for each singular value.
\end{proof}

These are conditional estimates on one target block with its parent blocks
held fixed; Section~\ref{subsec:stability} explains why they do not by
themselves control the global certificate.

\subsection{Stable and unstable propagation routes}
\label{subsec:routes}

Write $\cond$ for the ratio of the larger to the smaller nonzero singular
value.

\begin{corollary}[uniform local conditioning of the explicit route]
\label{cor:conditioning}
Every post-seed local projected interaction block in the sparse solenoidal
ordering of Section~\ref{subsec:solprop} satisfies the sharp bounds
\begin{equation}
\label{eq:sharpcond}
\min(\sigma_{\mathrm n}^{2},\sigma_{\mathrm t}^{2})\ge\frac1{10},
\qquad
\cond\le\sqrt{15} .
\end{equation}
\end{corollary}

\begin{proof}
The three bootstrap triads in \eqref{eq:boot1} have, up to interchanging $q$
and $r$, $(\lambda,\mu,\kappa,Z^{2})=(1,2,3,2)$, hence
$(\sigma_{\mathrm n}^{2},\sigma_{\mathrm t}^{2})=(3,1/3)$. The three triads in
\eqref{eq:boot2} have $(\lambda,\mu,\kappa,Z^{2})=(3,1,2,2)$, hence
$(\sigma_{\mathrm n}^{2},\sigma_{\mathrm t}^{2})=(8/3,4/3)$.

The shifted axis modes use $q=ne_{i}$, $r=e_{i}+e_{j}$ with $n\ge1$. Direct
substitution in \eqref{eq:projspec} gives
\begin{equation}
\label{eq:axisspec}
\sigma_{\mathrm n}^{2}=\frac{n^{2}+2}{2},
\qquad
\sigma_{\mathrm t}^{2}=\frac{(n^{2}-2)^{2}}{2(n^{2}+2n+2)} .
\end{equation}
The second expression is at least $1/10$ for integers $n\ge1$, with equality at
$n=1$ and value $1/5$ at $n=2$; for $n\ge3$, comparison with $1$ reduces to
$n(n^{3}-6n-4)>0$, which holds at $n=3$ and is strictly increasing thereafter.
Moreover
\[
15(n^{2}-2)^{2}-(n^{2}+2)(n^{2}+2n+2)
=(n-1)(n-2)(14n^{2}+40n+28)\ge0
\]
for $n\ge1$, so $\sigma_{\mathrm n}^{2}/\sigma_{\mathrm t}^{2}\le15$, the cases
$n=1,2$ being equalities of the displayed inequality.

The completion of the positive axes uses $q=ne_{i}+e_{j}$, $r=-e_{j}$ with
$n\ge2$, giving
\begin{equation}
\label{eq:axisspec2}
\sigma_{\mathrm n}^{2}=\frac{n^{2}(n^{2}+2)}{n^{2}+1},
\qquad
\sigma_{\mathrm t}^{2}=\frac{n^{4}}{n^{2}+1} .
\end{equation}
Here the smaller squared singular value is at least $16/5$ and the squared
condition number is $(n^{2}+2)/n^{2}\le3/2$.

Finally consider a general non-axis step \eqref{eq:generalstep}. It chooses a
unit vector $r$ and a parent $q$ with
\[
m=|q|^{2}\ge2,
\qquad
a=q\cdot r,
\qquad
d=m-a^{2}=|q\times r|^{2}\ge1,
\qquad
0\le a\le\sqrt m .
\]
Formula \eqref{eq:projspec} becomes
\begin{equation}
\label{eq:generalspec}
\sigma_{\mathrm n}^{2}=\frac{d(m+1)}m,
\qquad
\sigma_{\mathrm t}^{2}=\frac{d(m-1)^{2}}{m(m+1+2a)} .
\end{equation}
The normal value exceeds $1$. If $m=2$, integrality forces $a\in\{0,1\}$ and
the tangent value is $1/3$ or $1/10$, with squared condition number $9$ or
$15$. If $m\ge3$, then
\[
\sigma_{\mathrm t}^{2}
\ge\frac{(m-1)^{2}}{m(\sqrt m+1)^{2}}
=\Bigl(1-\frac1{\sqrt m}\Bigr)^{2}>\frac1{10},
\]
while
\[
\frac{\sigma_{\mathrm n}^{2}}{\sigma_{\mathrm t}^{2}}
=\frac{(m+1)(m+1+2a)}{(m-1)^{2}}
\le\frac{m+1}{(\sqrt m-1)^{2}}
\le4+2\sqrt3<15 .
\]
The penultimate expression is decreasing in $\sqrt m>1$, so its maximum for
$m\ge3$ occurs at $m=3$ and equals $4+2\sqrt3$. This proves
\eqref{eq:sharpcond}; both constants are attained in \eqref{eq:axisspec} at
$n=1$, equivalently in \eqref{eq:generalspec} at $m=2$, $a=1$.
\end{proof}

Unequal parent lengths alone do not imply uniform conditioning. For
\[
q_{n}=(n,0,0),
\qquad
r_{n}=(0,n,1),
\qquad
p_{n}=-(q_{n}+r_{n}),
\]
one has $Z^{2}=n^{2}(n^{2}+1)$, $\lambda=n^{2}$, $\mu=n^{2}+1$ and
$\kappa=2n^{2}+1$, so that \eqref{eq:projspec} yields
\begin{equation}
\label{eq:badfamily}
\sigma_{\mathrm n}^{2}=2n^{2}+1,
\qquad
\sigma_{\mathrm t}^{2}=(2n^{2}+1)^{-1},
\qquad
\cond=2n^{2}+1 .
\end{equation}
The parents are never of equal length, yet the conditioning degenerates
quadratically. Choosing the propagation route well is therefore a genuine
constraint, not a formality.

Finally, since
$|\lambda-\mu|/\sqrt{\lambda\mu}
=2|\sinh(\tfrac12\log(\mu/\lambda))|$,
taking the square root of the second formula in \eqref{eq:projspec} gives
\begin{equation}
\label{eq:sinh}
\sigma_{\mathrm t}
=\frac{2Z}{\sqrt\kappa}
\Bigl|\sinh\Bigl(\tfrac12\log\tfrac\mu\lambda\Bigr)\Bigr| .
\end{equation}
Observability of the tangential channel is thus governed by the logarithmic
separation of the two parent shells.

\subsection{A finite stability estimate}
\label{subsec:stability}

Fix the normalized frames of Section~\ref{sec:setting} and a normalization of
the selected real test fields. Let $\mathsf A_{\Lambda}$ be the resulting
\emph{global} certificate matrix, let $x_{R}$ be the real coordinate vector of
the complete symbol, and let $K=\ker\mathsf A_{\Lambda}$ be the appropriate
classified kernel.

\begin{proposition}[global stability, and its distinction from local
conditioning]
\label{prop:stability}
If $\gamma_{\Lambda}$ is the smallest positive singular value of
$\mathsf A_{\Lambda}$, then
\begin{equation}
\label{eq:stability}
\dist(x_{R},K)\le\gamma_{\Lambda}^{-1}\|\mathsf A_{\Lambda}x_{R}\|_{2} .
\end{equation}
In propagation order the rows at a new target have the schematic form
$y_{j}=L_{j}x_{j}+\sum_{\ell<j}E_{j\ell}x_{\ell}$, where $L_{j}$ is the local
diagonal block and the $E_{j\ell}$ carry the already encountered parent
columns. Full rank of the seed block and of every $L_{j}$ proves the global
rank by block triangularity. It does not bound $\gamma_{\Lambda}$ in terms of
$\min_{j}\sigma_{\min}(L_{j})$ alone, because the off-diagonal blocks can
accumulate earlier errors during forward substitution.
\end{proposition}

\begin{proof}
Restrict $\mathsf A_{\Lambda}$ to $K^{\perp}$; its least singular value there
is $\gamma_{\Lambda}$, so
$\|\mathsf A_{\Lambda}x\|_{2}\ge\gamma_{\Lambda}\|x\|_{2}$ for
$x\in K^{\perp}$. Applying this to the orthogonal projection of $x_{R}$ onto
$K^{\perp}$ gives \eqref{eq:stability}. The block formula follows from the
propagation ordering: each selected row involves its current target and only
earlier parent blocks. Triangularity therefore proves rank, whereas the global
least singular value depends on the complete matrix, including the
$E_{j\ell}$.
\end{proof}

Thus, for each fixed finite certificate, approximate cancellation on its
selected fields quantitatively forces proximity to the classified family. This
is a finite-certificate stability estimate and nothing more: the local bounds
\eqref{eq:sharpcond} imply neither an $N$-independent nor a polynomial bound
for $\gamma_{\Lambda}^{-1}$, and no such global bound is claimed here. Making
one is Problem~\ref{prob:gamma} below.

\section{Strain--vorticity applications}
\label{sec:strainvorticity}

This section transports the classification through the vorticity and records
its exact structural consequence for anisotropic second-order operators. It is
deliberately bounded: what follows are identities, not estimates.

\subsection{The curl pullback and the proof of Corollary~C}
\label{subsec:pullback}

Let $u$ be admissible, $\omega=\curl u$ and $S=\gsym u$, and let
$T(k):\kperp\to\C^{3}$ be complex linear and reality compatible. No
solenoidal-output assumption is imposed. With $\Dform_{T}$ as in
\eqref{eq:Dform}, symmetry of $\omega_{i}\omega_{j}$ gives
\[
\Dform_{T}(u)
=\int_{\T^{3}}\partial_{i}(Tu)_{j}\,\omega_{i}\omega_{j}\,dx ,
\]
and since $\partial_{i}\omega_{i}=0$, periodic integration by parts gives
\begin{equation}
\label{eq:byparts}
\Dform_{T}(u)
=-\int_{\T^{3}}(Tu)_{j}\,\omega_{i}\partial_{i}\omega_{j}\,dx
=-\langle Tu,(\omega\cdot\nabla)\omega\rangle .
\end{equation}
Note that $\Dform_{T}(u)$ is \emph{not} $\langle Tu,S(u)\omega\rangle$; that
is a different cubic form.

On mean-zero solenoidal fields the curl is a bijection: for $k\ne0$ and
$v\perp k$ one has $\Ck^{2}v=|k|^{2}v-k(k\cdot v)=|k|^{2}v$ and hence
$\Ck^{-1}v=\Ck v/|k|^{2}$, and this inverse preserves finite support, reality
and transversality. Setting
\begin{equation}
\label{eq:substitution}
w=\omega,
\qquad
R=T\curl^{-1} ,
\end{equation}
equation \eqref{eq:byparts} becomes the exact equivalence
\begin{equation}
\label{eq:equivalence}
\Dform_{T}(u)=-\Qform_{R}(w) .
\end{equation}
The minus sign and the single factor $|k|^{-2}$ are both essential. No growth
condition on either symbol is needed, because the tests are finite Fourier
sums.

\begin{proof}[Proof of Corollary~C]
As $u$ ranges over admissible fields, $w=\curl u$ ranges over all admissible
fields, so by \eqref{eq:equivalence} the cancellation $\Dform_{T}\equiv0$ is
equivalent to $\Qform_{R}\equiv0$ for $R=T\curl^{-1}$. Theorem~A gives
$R=A+B\curl+\curl B$ with unique constant real symmetric $A,B$. Composing on
the right with $\curl$ and using $\curl^{2}=-\Delta$ on the mean-zero
solenoidal space gives
\[
T=A\curl+B(-\Delta)+\curl B\curl ,
\]
which is \eqref{eq:Tfamily}; uniqueness is inherited from Theorem~A and the
invertibility of the curl. This is a bijective pullback of the
classification, not an extra condition forcing $A$ or $B$ to be scalar.

If $T$ has solenoidal output then so does $R=T\curl^{-1}$, and Corollary~B
gives $R=cI+d\curl$, hence $T=c\curl+d(-\Delta)$, which is \eqref{eq:Tsol}.
Conversely these operators preserve divergence-free output and satisfy the
cancellation. Formula \eqref{eq:Tsol} requires the range assumption;
\eqref{eq:Tfamily} does not.
\end{proof}

Taking $T=-\Delta$ in \eqref{eq:Tsol} recovers Miller's scalar Laplacian
identity~\cite{Miller2026}
\begin{equation}
\label{eq:miller}
\langle-\Delta S,\omega\otimes\omega\rangle=0 .
\end{equation}
It does \emph{not} give $\langle S,\omega\otimes\omega\rangle=0$: the field
\[
u=2(0,1,1)\cos x+2(1,0,0)\cos y-2(1,-1,1)\sin(x+y)
\]
is real, mean zero and divergence free, and satisfies
$\int\omega^{T}S\omega=-2$ in normalized measure. Corollary~C therefore says
that, within the reality-compatible translation-invariant multiplier class on
$\T^{3}$ and for the universal test family, no scalar, anisotropic,
matrix-valued, nonlocal or non-self-adjoint solenoidal Fourier multiplier
outside $\spn\{\curl,-\Delta\}$ creates a new universal strain--vorticity
cancellation.

For the quadratic strain energy
\begin{equation}
\label{eq:strainenergy}
\mathcal E_{T}(u)
=\tfrac12\langle S,\gsym(Tu)\rangle
=\tfrac14\langle-\Delta u,Tu\rangle ,
\end{equation}
call $T$ coercive when there is $c_{0}>0$ such that
$\mathcal E_{T}(u)\ge c_{0}\|u\|_{\dot H^{2}}^{2}$ for every admissible $u$.
Within the classified solenoidal family, the complete coercive cone on the
mean-zero $2\pi$-torus is
\begin{equation}
\label{eq:cone}
T=a(-\Delta)+b\curl,
\qquad
a>|b| .
\end{equation}
Indeed, on a helical Fourier mode the curl eigenvalue is
$\pm|k|$, so the contribution to \eqref{eq:strainenergy} is
\[
\frac14|k|^{2}\bigl(a|k|^{2}\pm b|k|\bigr)|\widehat u(k)|^{2}.
\]
Because every nonzero lattice frequency satisfies $|k|\ge1$, uniform
positivity relative to $|k|^{4}|\widehat u(k)|^{2}$ is equivalent to
$a>|b|$; necessity is already witnessed by the two helicities at
$|k|=1$. This is algebraic coercivity together with cancellation of one
nonlinear channel. It is not a closed PDE estimate.

\subsection{Anisotropic longitudinal completion}
\label{subsec:anisotropic}

The unrestricted family \eqref{eq:Tfamily} answers a question the solenoidal
one cannot: given an arbitrary constant-coefficient anisotropic scalar
second-order operator, what must be added to it to restore the universal
strain--vorticity cancellation?

For a real symmetric matrix $H$ set
\begin{equation}
\label{eq:Hcorrespondence}
B_{H}=\frac{\tr H}{2}I-H,
\qquad
L_{H}=-\sum_{i,j}H_{ij}\partial_{i}\partial_{j},
\qquad
\Phi_{H}(u)=\dv(B_{H}u).
\end{equation}
The correspondence is invertible, with inverse $H=(\tr B)I-B$; the two
formulas are different, and $H\mapsto B_{H}$ is not an involution on
$\Sym(3,\R)$.

\begin{corE}[anisotropic longitudinal completion]
\label{cor:E}
Let $H\in\Sym(3,\R)$ and put $T_{B}=B(-\Delta)+\curl B\curl$. Then on
divergence-free fields
\begin{equation}
\label{eq:bridge}
T_{B_{H}}u=L_{H}u-\nabla\Phi_{H}(u),
\qquad
PT_{B_{H}}u=L_{H}u ,
\end{equation}
and $B_{H}$ is the unique symmetric matrix with this property. Consequently
every constant-coefficient anisotropic scalar second-order operator has a
unique homogeneous second-order completion of the form $T_{B}$ inside the
classified family, and the completed operator
satisfies
\begin{equation}
\label{eq:completed}
\bigl\langle
\gsym\bigl(L_{H}u-\nabla\Phi_{H}(u)\bigr),\omega\otimes\omega
\bigr\rangle=0 ,
\end{equation}
and the defect of the bare anisotropic operator is exactly longitudinal:
\begin{equation}
\label{eq:defect}
\langle L_{H}S,\omega\otimes\omega\rangle
=\langle\nabla^{2}\Phi_{H}(u),\omega\otimes\omega\rangle .
\end{equation}
\end{corE}

\begin{proof}
A coordinate expansion for symmetric $B$ gives, for \emph{every} vector $v$,
\begin{equation}
\label{eq:symbolidentity}
(B\Ck+\Ck B)\Ck v=(k^{T}Hk)v-(Hk\cdot v)k,
\qquad H=(\tr B)I-B .
\end{equation}
On $v\perp k$ the left-hand side is $|k|^{2}Bv+\Ck B\Ck v$, which is the
symbol of $T_{B}$. One must not replace $\Ck B\Ck v$ by
$|k|^{2}Bv-k(k\cdot Bv)$: that purported matrix triple-product identity is
false. For instance with $H=\mathrm{diag}(1,2,3)$, so
$B_{H}=\mathrm{diag}(2,1,0)$, and $k=e_{1}$, $v=e_{2}$, one has
$\Ck B_{H}\Ck v=0$. Hence the actual full left-hand side of
\eqref{eq:symbolidentity} is
$B_{H}\Ck^{2}v+\Ck B_{H}\Ck v=e_{2}$. Replacing only the second summand by
the incorrect expression would instead make the full left-hand side equal to
$e_{2}+e_{2}=2e_{2}$, while the right-hand side of
\eqref{eq:symbolidentity} is $e_{2}$. Thus the proposed replacement fails
when like quantities are compared.

For transverse $v$, symmetry and $B_{H}=(\tr H/2)I-H$ give
$k\cdot B_{H}v=-Hk\cdot v$, so \eqref{eq:symbolidentity} reads
$T_{B_{H}}u=L_{H}u-\nabla\Phi_{H}(u)$ in physical space, and applying the
Leray projection $P$ kills the gradient. This is \eqref{eq:bridge}.

For uniqueness, suppose two completions have the same projection and let $G$
be the difference of the corresponding $H$ matrices. Formula
\eqref{eq:symbolidentity} gives $(k^{T}Gk)v=0$ for every nonzero lattice $k$
and every $v\perp k$. Taking $k=e_{i}$ forces the diagonal of $G$ to vanish and
taking $k=e_{i}+e_{j}$ forces its off-diagonal entries to vanish, so $G=0$; the
inverse correspondence then gives equality of the two $B$ matrices.

Identity \eqref{eq:completed} is Corollary~C applied to
$T=T_{B_{H}}$, which is the case $A=0$, $B=B_{H}$ of \eqref{eq:Tfamily}; it is
not a direct application of Euler cancellation to $T_{B}$. Applying the
symmetric gradient to \eqref{eq:bridge} and substituting into
\eqref{eq:completed} gives \eqref{eq:defect}.
\end{proof}

The longitudinal term cannot be discarded from the pairing in
\eqref{eq:defect}, because $(\omega\cdot\nabla)\omega$ need not be divergence
free; an argument pairing a longitudinal vector against a transverse one is
therefore insufficient to eliminate this defect.

Since $\dv u=0$ and $H$ is symmetric, $\Phi_{H}(u)=-H_{ab}S_{ab}$. Expanding
\eqref{eq:defect} over all symmetric $H$ gives the integrated tensor identity
\begin{equation}
\label{eq:tensoridentity}
\int_{\T^{3}}\omega_{i}\omega_{j}
\bigl(\partial_{a}\partial_{b}S_{ij}-\partial_{i}\partial_{j}S_{ab}\bigr)\,dx=0,
\qquad 1\le a,b\le3 .
\end{equation}
The tensor is symmetric in $a,b$. Its trace is Miller's scalar identity
\eqref{eq:miller}, up to the immaterial overall sign --- it is not three times
that identity --- and its trace-free part supplies five anisotropic relations.
For $H=cI$ one has $\Phi_{H}=0$ and \eqref{eq:defect} reduces to
\eqref{eq:miller}.

\subsection{What does not follow}
\label{subsec:not-invariants}

Two boundaries must be stated explicitly.

\emph{The twelve directions are cancellations, not invariants.} When $R$ is
self-adjoint on the solenoidal phase space, $\Qform_{R}=0$ is the statement
that $\tfrac12\langle w,Rw\rangle$ is conserved by Euler. For general
$R:\kperp\to\C^{3}$ the longitudinal output pairs against the pressure
gradient in the full Euler equation, so the ten directions beyond
$\spn\{I,\curl\}$ are not, in general, new Euler invariants. The correct
reading of Theorem~A is a rigidity statement about the convective term, and
that is how it is used in Corollary~C.

\emph{The anisotropic completion is not a regularity criterion.} Identity
\eqref{eq:defect} is exact, and for fixed $H$ the elementary bound
\[
|\langle L_{H}S,\omega\otimes\omega\rangle|
\le\|\omega\|_{2}^{2}\,\|\nabla^{2}\Phi_{H}(u)\|_{L^{\infty}(\mathrm{op})}
\]
holds; it bounds the $H$-contracted expression, not every individual component
for a single choice of $H$. The defect is a constant-coefficient third
derivative of $u$, and standard energy control supplies no bound for the
displayed norm. It is moreover a pressure-\emph{like} longitudinal potential
rather than a nonlinear pressure: $\Phi_{H}$ is linear in $u$, whereas a
pressure term is determined quadratically, and we assert no formula identifying
their Hessians. For nonscalar $H$ the completion does not remove the analytic
difficulty; it relocates it into the Hessian channel. No new a priori,
scale-critical, time-integrable bound for that defect is proved here, and
the geometric-depletion program is context and a possible direction, not a
theorem of this paper.

\section{Computer-assisted components, reproduction and formalization}
\label{sec:computer}

\subsection{What is analytic and what is finite}
\label{subsec:split}

The proof is not a numerical experiment. Its logical division is as follows.
\begin{itemize}
\item \emph{Analytic}: sufficiency of the generators
(Section~\ref{sec:sufficiency}); phase separation (Lemma~\ref{lem:phase}); the
raw and projected triad lemmas (Lemmas~\ref{lem:rawline} and
\ref{lem:projected}); propagation through $\Z^{3}$
(Section~\ref{sec:propagation}); reduction by the range condition
(Corollary~B); uniqueness of $A,B$ (Proposition~\ref{prop:uniqueness}); the
curl pullback and the anisotropic completion
(Section~\ref{sec:strainvorticity}); and the singular-value formulas
(Section~\ref{sec:singular}).
\item \emph{Computer-assisted but exact}: the ranks of the two finite seeds,
the selection of explicit nonsingular minors, the independent reconstruction
of the constraint matrices, the unit-cube minimality search, and a large finite
replay of the propagation order.
\item \emph{Finite-dimensional linear algebra}: the optimality and stability of
the sparse certificates (Theorem~D and
Proposition~\ref{prop:stability}).
\end{itemize}
No floating-point rank decision and no numerical tolerance enters the
certificates. The computations use integers, exact rational arithmetic, or
exact reductions modulo primes.

\subsection{Reproduction}
\label{subsec:repro}

The exact programs are run from a single entry point. They execute in isolated
subprocesses and report the interpreter version, the source revision, the
working-tree state and the elapsed time; the run fails if any certificate is
missing or any assertion fails. The checks cover
\begin{itemize}
\item the $720$ solenoidal seed rows, exact rank $54/56$, and the two minor
residues \eqref{eq:solminor};
\item an independent replay of the same rows and of the energy--helicity
kernel;
\item all $264$ signed non-collinear unit-box target triads, with the exact
complex ranks $8/4$ of \eqref{eq:localranks};
\item $5{,}199$ solenoidal propagation steps, cumulatively over the boxes of
radii $2,3,4,10$ (the radius-ten replay alone has $4{,}623$ steps);
\item the $680$ unrestricted-output rows, exact modular rank $144/156$, the
twelve generators, and the two minor residues \eqref{eq:rawminor};
\item an independent reproduction of the unrestricted rows and of their
SHA-256 digest;
\item $4{,}617$ raw two-triad target blocks at exact rank twelve, reaching all
$9{,}260$ signed modes through radius ten;
\item all $4{,}095$ unit-cube subsets of at most six modes in the seven-mode
minimality audit;
\item sharp Hermitian certificates through radii $1,2,3,4,10$, including the
assembled rational rank $246=4\cdot62-2$ at radius two;
\item an exact semantic reconstruction of all $26$ selected Hermitian
Boolean-seed rows from their signed triads, transverse polarizations and phase
choices, checked row-by-row against the typed Lean matrix up to nonzero
rational scaling;
\item $56{,}064$ signed non-collinear triads, cumulatively over the boxes of
radii $1,2,3$, against the exact singular-value formulas
\eqref{eq:rawspec} and \eqref{eq:projspec};
\item the sharp local minimum $1/10$, the sharp maximum condition number
$\sqrt{15}$ of \eqref{eq:sharpcond}, and the near-isosceles family
\eqref{eq:badfamily} through $n=100$;
\item the curl-pullback identity \eqref{eq:equivalence}, the explicit negative
control after \eqref{eq:miller}, and the corrected cancellations, in rational
Gaussian arithmetic.
\end{itemize}
The selected seed rows, pivot columns, mode order and parameter order are
exported in the accompanying machine-readable file
\texttt{certificate-index.json} in the supplement directory. Row digests and
modular-minor residues are asserted inside the individual programs and
summarized in the repository certificate manifest. Two of the
programs are deliberately independent at the algorithmic level: they import
neither the corresponding Fourier row builders nor the row reducers, and
instead reconstruct the convolution by recipient wavevector. Their agreement
therefore checks the row construction through a distinct implementation path. Both are nevertheless
written in the same language, so a re-implementation in a second system would
strengthen the evidence further.

\subsection{Formalization and the trust boundary}
\label{subsec:lean}

Theorem~A, its all-mode completion, Corollaries~B and C, the matrix-level
finite-dimensional content supporting Theorem~D, and the local spectral and
conditioning formulas are formalized in Lean~4 with Mathlib. The four
supporting results added after the original classification have the following
precise scope.
\begin{itemize}
\item \emph{Theorem~A$^{0}$.} Lean classifies every nonzero transverse
fiber, identifies reality compatibility at frequency zero with a real
$3\times3$ block, and proves that replacing this block by any real matrix
does not change universal all-mode cancellation.
\item \emph{Theorem~D.} Lean proves the rank--nullity lower bound for real
scalar evaluation matrices, constructs sharp abstract coordinate
certificates with counts $8M-2$, $4M-2$ and $12M-12$, verifies the
$26\times28$ Hermitian Boolean-seed matrix has nullity two, and checks the
ordered-count identities. Lean also proves a general propagation theorem
for arbitrary-length unequal-triad orderings of any mode-set size, instantiated
at one additional mode beyond the Boolean seed with explicit named test
fields whose Euler cubic pairing vanishes against every admissible symbol.
The realization of the complete sharp geometric families for every admissible
ordering remains the analytic and exact-program argument of
Sections~\ref{sec:results} and~\ref{sec:propagation}.
\item \emph{Corollary~C.} For trigonometric polynomials, Lean evaluates the
velocity and its vorticity as functions on the torus, differentiates them with
the Mathlib derivative rather than by symbol manipulation, forms the pointwise
contraction $\gsym(Tu):(\omega\otimes\omega)$, and proves that its normalized
integral over $\mathbb{T}^{3}$---an iterated interval integral over one full
period in each coordinate, reduced by the character orthogonality relation
$\int_{0}^{2\pi}e^{\mathrm{i}nt}\,\mathrm{d}t=0$ for $n\neq0$---equals the
finite-Fourier functional, together with the physical-space identity
$\int\gsym(Tu):(\omega\otimes\omega)=-\int Tu\cdot(\omega\cdot\nabla)\omega$.
The classification transfers to that integral. A second layer extends the
functional by an absolutely convergent resonant series to the smooth
solenoidal fields, those whose Fourier coefficients decay faster than every
polynomial, exhibits a member of that class with infinite spectrum, and proves
the same classification there by truncation; the extended functional is tied
to physical space as the limit of the torus integrals of its truncations.
No Sobolev completion at finite regularity is asserted.
\item \emph{Section~\ref{sec:singular}.} Lean derives the raw and projected
Gram matrices, the multiplicity-four observation spectrum, the sharp $1/10$
and $\sqrt{15}$ bounds for every route family---the shifted-axis steps, the
general steps with $m\ge3$ and the two admissible $m=2$ steps, and the two
bootstrap triads, with the axis formulas obtained by substituting the
axis-step parameters into the projected Gram spectrum rather than postulated
separately---the quadratically ill-conditioned family, and the stability
estimate \eqref{eq:stability} with no spectral hypothesis: positivity of
$\gamma_{\Lambda}$ on $K^{\perp}$ is derived from finite-dimensionality
rather than assumed, and the estimate is instantiated at the Hermitian
Boolean-seed certificate itself. The constant is per certificate; no bound
uniform in the mode count is formalized, in line with
Problem~\ref{prob:gamma}.
\end{itemize}
At the checked revision, the source contains no \texttt{sorry} or
\texttt{admit}, and every declared module elaborates under the pinned
toolchain. This does not make every certificate kernel-checked; the trust
status splits as follows.

Kernel-checked, depending only on \texttt{propext},
\texttt{Classical.choice} and \texttt{Quot.sound}: the vector identities used
throughout; Lemma~\ref{lem:scalarity} in a slightly stronger form that does not
assume symmetry; Proposition~\ref{prop:uniqueness} and the uniqueness of the
two generator matrices on all nonzero transverse fibers; the finite-support
Fourier definitions of admissible fields, symbols, convolution and the cubic
form; reality compatibility of the classified family; the generator
cancellations on every triad and the finite-sum argument connecting triad
cancellation to the convolutional cubic form, that is, sufficiency for every
admissible finite Fourier field including an arbitrary invisible zero-mode
block; phase polarization, i.e.\ Lemma~\ref{lem:phase}; the raw-recipient
parametrization and the local two-triad determination rule of
Lemma~\ref{lem:rawline}; reality transport of fiberwise vanishing; the
propagation theorem in the $\ell^{1}$-shell order of
Section~\ref{subsec:rawprop}, with concrete parent-pair certificates for every
nonzero lattice mode outside the cube; the two-frame decomposition of a
transverse fiber at each of the thirteen seed modes; a non-vacuity witness
showing that $R(k)=\|k\|^{2}I$ is reality compatible but not a generator, so
that neither side of the classification is trivially satisfied; and the
\emph{semantic seed bridge}, which proves that a reality-compatible symbol with
triad cancellation satisfies every selected constraint of the unit-cube matrix.
The bridge is what makes the rank computation a statement about symbols rather
than about an arbitrary integer matrix.

Compiler-trusted: the rank-$144$ claim over $\Z/101\Z$ and the explicit inverse
certificate for the selected $144\times144$ modular minor are discharged by
\texttt{native\_decide}, which evaluates compiled code instead of reducing in
the kernel, and therefore trusts the Lean compiler, its runtime and GMP. The
formal statements of Theorem~A and Corollary~B inherit seven such axioms. The
precise status is: \emph{the analytic necessity and sufficiency arguments
for Theorem~A and Corollary~B are kernel-checked, and the formal
Theorem~A holds modulo trusting the Lean compiler for one finite
linear-algebra computation.} The certificate is roughly three million modular
operations, which is not viable for kernel reduction at present; shrinking it
is the route to removing this dependency. The authoritative check is
\texttt{\#print axioms} applied to the formal statements, not any prose list.

The later formalization has three deliberately separated trust layers. First,
the all-mode theorem records the arbitrary real zero block; Corollary~C is
realized as a normalized integral over the three-torus of the pointwise
contraction, and extended from trigonometric polynomials to the smooth
solenoidal fields; the geometric frame identities, local Gram spectra,
multiplicity-four observation, route bounds, bad family and stability modulo
the exact kernel are kernel-checked. Second, Lean proves that the typed
$26\times28$ Hermitian Boolean-seed matrix has rank $26$ and nullity $2$ and
proves the sharp $8M-2$, $4M-2$ and $12M-12$ scalar lower bounds at the
finite-dimensional matrix level; the seed rank calculation uses the disclosed
\texttt{native\_decide} boundary. Third, the mandatory exact program
\texttt{verify\_solenoidal\_seed\_semantics.py} reconstructs every selected
Hermitian row from an explicit triad-supported Euler test field and compares
it with the Lean matrix. This last physical-to-matrix reconstruction is an
exact executable bridge, not yet an internal kernel-checked Lean theorem.

The formal development stops at the smooth solenoidal fields. It does not
assert a Sobolev completion at finite regularity, a multiplier-domain theorem
or PDE solution theory, and on the smooth class the functional is tied to
physical space as the limit of the torus integrals of its truncations rather
than by an integral of the smooth field itself.
Likewise, the Lean matrix-count theorem does not by itself construct the full
geometric observation family for every ordering in Theorem~D; that construction
is the analytic and exact-certificate argument in Sections~\ref{sec:results}
and~\ref{sec:propagation}. These boundaries prevent an abstract coordinate
matrix from being mislabeled as a fully internalized physical certificate.

\section{Discussion and open problems}
\label{sec:discussion}

Universal orthogonality to three-dimensional Euler convection is highly
restrictive. Translation invariance and reality alone force an arbitrary Fourier symbol into a
finite-dimensional family of restrictions of constant-coefficient, formally
self-adjoint operators: twelve dimensions before the pressure direction is
removed and only two on the solenoidal phase space. Allowing longitudinal
output creates the larger kernel, while imposing the solenoidal range
condition leaves only the energy--helicity family. The isosceles defect is
created by the Leray projection and is measured by \eqref{eq:sinh}; both
statements admit sharp finite certificates
verified in exact arithmetic. Through the vorticity, the same theorem
classifies the universal Fourier-multiplier extensions of the scalar
strain--vorticity cancellation and identifies the precise pressure-like
completion that anisotropy requires.

The claims have clear limits. The paper does not address existence,
regularity or blow-up. It does not prove a multiplier theorem on $\R^{3}$, and
the zero Fourier block is invisible
rather than rigid. The novelty statement in Section~\ref{subsec:boundary}
is subject to the same limits.

Open problems include the following.
\begin{enumerate}
\item Can the finite seed ranks be replaced by a short symbolic argument that
preserves the unrestricted hypotheses, removing the computer-assisted step
entirely?
\item Does a six-mode solenoidal seed exist outside the coordinate unit cube?
The minimality result of Section~\ref{subsec:solseed} is confined to that cube.
\item\label{prob:gamma} Can the smallest positive singular value
$\gamma_{\Lambda}$ of the complete sparse certificate be bounded uniformly, or
at least polynomially, in the box radius? Proposition~\ref{prop:stability}
shows that the local bounds \eqref{eq:sharpcond} do not suffice.
\item What is the optimal propagation graph for global conditioning, as
opposed to the optimal number of tests?
\item Is there an analogous classification on $\R^{3}$ under natural
regularity assumptions on a continuous symbol?
\item How does the classification change on other compact three-manifolds, or
on irrational tori?
\item Can the logarithmic shell gap in \eqref{eq:sinh} be coupled to a
dynamical inter-shell decorrelation estimate?
\item Can the longitudinal completion \eqref{eq:defect} be exploited
dynamically without losing control to the Hessian channel?
\end{enumerate}

\appendix

\section{Certificate data}
\label{app:certificates}

The following constants are asserted by the verification programs and
summarized in the machine-readable manifest that accompanies the source
release.

\begin{center}
\small
\begin{tabular}{ll}
\hline
Quantity & Value\\
\hline
Complex solenoidal seed & $7$ unoriented modes, $56$ parameters\\
\quad nonzero rows / rank / nullity & $720$ / $54$ / $2$\\
\quad minor residues mod $(1000003,\,1000000007)$
 & $(190194,\,891695525)$\\
Hermitian solenoidal seed & $7$ unoriented modes, $28$ parameters\\
\quad nonzero rows / rank / nullity & $592$ / $26$ / $2$\\
Unrestricted-output seed & $13$ unoriented modes, $156$ parameters\\
\quad nonzero rows / rank / nullity & $680$ / $144$ / $12$\\
\quad minor residues mod $(1000003,\,1000000007)$
 & $(272144,\,290707620)$\\
\quad canonical row digest (SHA-256)
 & \texttt{a4fc3b45\dots39e9c30d}\\
Local projected ranks in the unit box
 & $264$ signed triads: $192$ of rank $8$, $72$ of rank $4$\\
Solenoidal propagation replay
 & $5199$ steps over radii $2,3,4,10$\\
Raw two-triad propagation replay
 & $4617$ target blocks, $9260$ signed modes\\
Unit-cube minimality search & $4095$ subsets, exactly four minimizers\\
Local conditioning replay & $56064$ signed triads over radii $1,2,3$\\
Sharp local bounds
 & $\sigma_{\min}^{2}=1/10$, $\cond=\sqrt{15}$\\
\hline
\end{tabular}
\normalsize
\end{center}

The full digest is
\begin{center}
\texttt{a4fc3b45a6fc94defd38d992dbbe27f4413254b3af307e41d1c7e25839e9c30d}
\end{center}
The solenoidal propagation count is cumulative over the listed boxes; the
radius-ten ordering alone contributes $4623$ steps. The local-conditioning
count is likewise cumulative. The unit-cube minimality statement exhausts all
subsets of at most six modes of $\Lambda_{1}$ and finds exactly four
seven-mode minimizers; it is not a statement about arbitrary six-mode
configurations elsewhere in the lattice.

\section*{Data and code availability}

The exact-arithmetic certificate programs, the machine-readable certificate
manifest and the Lean formalization accompany this paper in the
\href{https://github.com/raelnogpires/rigidity-of-fourier-multipliers-for-universal-euler-convective-cancellations}{GitHub repository}.
The materials corresponding to this version are archived under release
\href{https://github.com/raelnogpires/rigidity-of-fourier-multipliers-for-universal-euler-convective-cancellations/releases/tag/v1.0.0}{\texttt{v1.0.0}}.
The verification entry point and software versions are recorded in the
accompanying reproduction documentation.

\bibliographystyle{plain}
\bibliography{references}

\end{document}
